% concatenated from regsing.tex
\documentclass[a4paper, 11pt]{article}

\usepackage{TSeng}
\addbibresource{diffgaloisrefs.bib}

\title{The regular singular inverse problem in differential Galois theory}
\author{Thomas Serafini and Michael Wibmer\thanks{Supported by the Lise Meitner grant M-2582-N32 of the Austrian Science Fund FWF.}}
\date{\today}

\begin{document}

\maketitle

\begin{abstract}
	We show that every linear algebraic group over an algebraically closed field of characteristic zero is the differential Galois group of a regular singular linear differential equation with rational function coefficients. %use a recent result on specialization of differential Galois groups and differential torsors from \cite{fengwib}, along with the Riemann-Hilbert correspondence on $\C$ to realize any algebraic group over any algebraically closed field of characteristic zero as the Galois group of a regular singular differential equation.
	\let\thefootnote\relax\footnotetext{{\em Mathematics Subject Classification Codes:} 34M50, 12H05.
		{\em Key words and phrases}:
		Differential Galois theory, regular singular differential equations, Riemann-Hilbert correspondence, free proalgebraic groups.
		%		Michael Wibmer, Institute of Analysis and Number Theory, Graz University of Technology, Kopernikusgasse 24, 8010 Graz, Austria, \texttt{wibmer@math.tugraz.at}
	}
\end{abstract}

%??include comment from Michael Singer??

	
\section{Introduction}
Hilbert's 21st problem, also known as the Riemann-Hilbert problem, asks whether there exists a Fuchsian linear differential equation on the Riemann sphere with prescribed singularities and monodromy group. In general, the answer is no, but J. Plemelj showed that the answer is yes if the Fuchsian requirement is slightly relaxed by allowing one singularity to be regular singular (\cite{AnasovBolibruch:TheRiemannHilbertProblem}). In this article, we generalize Plemelj's result from the field of complex numbers to an arbitrary algebraically closed field of characteristic zero. As the monodromy group is only defined over the complex numbers, it is replaced with the differential Galois group in our context. This replacement is justified by Schesinger's density theorem, stating that, for regular singular linear differential equations, the differential Galois group is the Zariski closure of the monodromy group.

\begin{theorem}[Theorem \ref{thm:inv}] \label{theo: main}
	%	Let $G\leq\GL_n(k)$ be a linear algebraic group over $k$ and let $d\geq 0$ be such that $G$ is the Zariski closure of a subgroup generated by $d$ elements. Then, for every subset $X$ of $k$ consisting of $d$ elements, there exists a regular singular differential equation over $k(x)$ of order $n$ with singularities contained in $X\cup\{\infty\}$  and differential Galois group $G$.
	Let $k$ be an algebraically closed field of characteristic zero and let $S=\{p_1,\ldots,p_{d+1}\}$ be a subset of $\P^1(k)$ with $d+1$ elements. Then, for every Zariski closed subgroup $G$ of $\GL_{n,k}$ that can be generated by $d$ elements (as an algebraic group), there exists an ${n\times n}$-matrix $A$ of rational functions over $k$ such that the differential equation $\del(y)=Ay$ has differential Galois group $G$. Furthermore, $A$ can be chosen such that $\del(y)=Ay$ is Fuchsian at $p_1,\ldots,p_d$, regular singular at $p_{d+1}$ and all poles are contained in $S$.	
	 %In particular, every linear algebraic group can be realized as the differential Galois group of a regular singular differential equation.
\end{theorem}
Recall that $\del(y)=Ay$ is \emph{Fuchsian} at $p\in k$ if $A$ has at most a pole of order one at $p$ and that $\del(y)=Ay$ is \emph{regular singular} at $p$ if $\del(y)=Ay$ is Gauge-equivalent to an equation that is Fuchsian at $p$. %Furthermore, $p$ is a pole of $\del(y)=Ay$ if $p$ is a pole of one of the entries of $A$. 
These definitions extend to the point $p=\infty$ via the coordinate transformation $z\mapsto \frac{1}{z}$. The differential equation $\del(y)=Ay$ is \emph{regular singular} if it is regular singular at every point of $\P^1(k)$.

We note that in Theorem \ref{theo: main} for $k=\mathbb{C}$, the topological fundamental group of $\P^1(\C)\smallsetminus S$ is the free group on $d$ generators. Therefore, the monodromy group of $\del(y)=Ay$ can be generated by $d$ elements (as an abstract group) and the differential Galois of $\del(y)=Ay$ can be generated by $d$ elements (as an algebraic group). Furthermore, for $k=\mathbb{C}$, Theorem \ref{theo: main} is well-known. It is essentially the combination of Plemelj's (weak) solution of the Riemann-Hilbert problem with Schlesinger's density theorem. 

As every linear algebraic group over $k$ is the Zariski closure of a finitely generated subgroup and can be embedded into some general linear group, we obtain from Theorem \ref{theo: main}:

\begin{corollary}[Corollary \ref{cor:inv}] \label{cor:intro}
Every linear algebraic group is the differential Galois group of a regular singular differential equation. 
\end{corollary}	
Corollary \ref{cor:intro} generalizes \cite[Theorem 4.3]{Singer:ModuliOfLinearDifferentialEquationsOnTheRiemannSphereWithFixedGaloisGroups} where it was shown, based on a study of moduli of (scalar) linear differential equations, that Corollary \ref{cor:intro}
holds for a large class of linear algebraic groups, including connected and finite linear algebraic groups and groups whose identity component is semisimple or unipotent.


Corollary \ref{cor:intro} can be seen as a sharpening of the solution of the inverse problem in differential Galois theory, which states that every linear algebraic group is the differential Galois group of a linear differential equation with rational function coefficients (\cite{inverse}). Interestingly, this was first proved for $k=\C$ in \cite{TretkoffTretkoff:SolutionOfTheInverseProblem} via Theorem \ref{theo: main}, i.e., by combining Plemelj's (weak) solution of the Riemann-Hilbert problem and Schlesinger's density theorem. 


%Moreover, in the case $k=\mathbb{C}$, our proof of Theorem~\ref{theo: main} essentially reduces to the argument of \cite{TretkoffTretkoff:SolutionOfTheInverseProblem}. 
Our strategy to prove Theorem \ref{theo: main} is to deduce it from the case $k=\C$ via a specialization argument, based on a specialization theorem from \cite{fengwib}. This specialization theorem was used in \cite{fengwib} to generalize the solution of the inverse problem in a different direction, by showing that the absolute differential Galois group of the rational function field is a free proalgebraic group.
%Building on earlier work by several authors, the solution of the inverse problem over arbitrary (algebraically closed, characteristic zero) base fields was finally completed in \cite{inverse}.

%Let $k$ be an algebraically closed field of characteristic zero. We consider linear differential equations over $k(z)$, the field of rational functions with coefficients in $k$, equipped with the usual derivation $\frac{d}{dz}$. The \emph{solution of the inverse problem} in differential Galois theory states that every linear algebraic group over $k$ is the differential Galois group of a linear differential equation with coefficients in $k(z)$. This was first proved in \cite{inverse}, based on previous work of several authors.
%
%There are various directions into which the inverse problem can be generalized. For example, instead of trying to realize linear algebraic groups as differential Galois groups, one can try to realize surjective morphisms of linear algebraic groups. This idea can be formalized through differential embedding problems: Recall that a \emph{differential embedding problem} for $k(z)$ is given by a Picard-Vessiot extension $L$ of $k(z)$ with differential Galois group $H$ together with a surjective morphism $G\to H$ of linear algebraic groups. A \emph{solution} consists of a Picard-Vessiot extension $M$ of $k(z)$, containing $L$, with differential Galois group $G$ such that the restriction map $G\to H$ of differential Galois groups corresponds to the given morphism $G\to H$. Note that the special case $H=1$ recovers the inverse problem. Building on previous work (\cite{BachmayrHarbaterHartmannWibmer:DifferentialEmbeddingProblems}, \cite{BachmayrHarbaterHartmannPop:LargeFieldsInDifferentialGaloisTheory}, \cite{BachmayrHarbaterHartmannWibmer:FreeDifferentialGaloisGroups}, \cite{BachmayrHarbaterHartmannWibmer:TheDifferentialGaloisGroupOfRationalFunctionField}), it was recently shown in \cite{fengwib} that every differential embedding problem for $k(z)$ has a solution.

%As the linear algebraic groups that can be realized as differential Galois groups are exactly the algebraic quotients of the absolute differential Galois group, determining the absolute differential Galois group of $k(z)$, can also be seen as a generalization of the inverse problem. Matzat's conjecture states that this group is the free proalgebraic group on $|k|$ generators. This 
%
%
%
%
%\begin{enumerate}
%	\item In how many essentially different ways does a given linear algebraic group $G$ occur as a differential Galois group? That is, what is the cardinality of the set of isomorphism classes of Picard-Vessiot extensions with differential Galois group isomorphic to $G$?
%	\item What is the absolute differential Galois group of $k(x)$? (The differential Galois groups are then the algebraic quotients of the absolute differential Galois group.) 
%\end{enumerate}
%
%These questions have recently been answered (\cite{BachmayrHarbaterHartmannWibmer:FreeDifferentialGaloisGroups},\cite{BachmayrHarbaterHartmannWibmer:TheDifferentialGaloisGroupOfRationalFunctionField},\cite{Wibmer:SubgroupsOfFreeProalgebraicGroupsAndMatzatsconjecture},\cite{FengWibmer:DifferentialGaloisGroupsSpecializationsAndMatzatsConjecture}): For every non-trivial algebraic group $G$, the cardinality in (i) is equal to $|k|$, the cardinality of $k$. The absolute differential Galois group of $k(x)$ is the free proalgebraic group on a set of cardinality $|k|$.

%Another way to sharpen the inverse problem is to restrict the class of differential equations via which one hopes to realize a given linear algebraic group as a differential Galois group. As the class of \emph{regular singular differential equations} is one of the best studied classes of linear differential equations, it is natural to wonder ``Is every linear algebraic group the Galois group of a regular singular linear differential equation?'' Our main result is that the answer is yes. In fact, we offer a somewhat more precise statement.
%
%\begin{theononumber}[Theorem \ref{thm:inv}]
%%	Let $G\leq\GL_n(k)$ be a linear algebraic group over $k$ and let $d\geq 0$ be such that $G$ is the Zariski closure of a subgroup generated by $d$ elements. Then, for every subset $X$ of $k$ consisting of $d$ elements, there exists a regular singular differential equation over $k(x)$ of order $n$ with singularities contained in $X\cup\{\infty\}$  and differential Galois group $G$.
%	Let $G\leq\GL_{n,k}$ be a linear algebraic group over $k$ that can be generated by $d$ elements (as an algebraic group) and let $S$ be a subset of $\P^1(k)$ with $d+1$ elements. Then there exists a matrix $A\in k(z)^{n\times n}$ such that the differential equation $\del(y)=Ay$ is regular singular with singularities contained in $S$ and has differential Galois group $G$. %In particular, every linear algebraic group can be realized as the differential Galois group of a regular singular differential equation.
%\end{theononumber}

%
%
%The above theorem is reminiscent of Schlesinger's density theorem: For the field $k=\C$ of complex numbers, the differential Galois group of a regular singular differential equation with $d+1$ singularities is the Zariski closure of the monodromy group, which is generated by $d$ elements.
%Moreover, the weak solution of the Riemann-Hilbert problem states that every subgroup of $\GL_n(\C)$ that can be generated by $d$ elements, occurs as the monodromy group of a regular singular differential equation over $\C(z)$ of order $n$ with the $d+1$ (potential) singularities being fixed a priori. Therefore, the above theorem holds for $k=\C$. This is well known and, in fact, this is exactly how the inverse problem was first solved over $\C(z)$ in \cite{TretkoffTretkoff:SolutionOfTheInverseProblem}. 
%

  

For $k=\mathbb{C}$ it is an immediate consequence of the Riemann-Hilbert correspondence that the differential Galois group of the family of all regular singular differential equations with singularities contained in $S$, is the free proalgebraic group on $d$ generators (where $|S|=d+1$). It seems natural to expect that this statement is true in full generality, i.e., for any algebraically closed field $k$ of characteristic zero. This question was already raised in \cite{Wibmer:RegSing} and the above theorem can be seen as a small step towards its resolution. Our theorem implies that every linear algebraic group over $k$ that is a quotient of the free proalgebraic group on $d$ generators also is a quotient of the differential Galois group of the family of all regular singular differential equations with singularities in $S$ (Corollary \ref{cor:proalg}).

\medskip

The authors are grateful to Michael Singer and the anonymous referee for helpful comments.

%For a subset $S$ of $\P^1(k)$ with $d$ elements it seems natural to expect that the differential  $k$ and let $\Gamma_X$ be the differential Galois group of the family of all regular singular differential equations over $k(x)$ with singularities contained in $S$. In \cite{Wibmer:RegularSingularDifferentialEquationsAndFreeProalgebraicGroups}, the question was raised if $\Gamma_X$ is the free proalgebraic group on $X$. Again, it is an immediate consequence of the Riemann-Hilbert correspondence that this is true for $k=\C$. Although we are far from resolving this question, Theorem A can be seen as a step in the right direction. It shows that every linear algebraic group that is a quotient of the free proalgebraic group on $X$, also is a quotient of $\Gamma_X$.


%The inverse problem in differential Galois theory consists of asking, given a differential field $(K, \del)$ and an algebraic group $G$ over  the field of constants $k = K^\del$ of $K$, whether there is a differential equation $\del y = Ay$ over $K$ with Galois group $G$.
%
%\newpar
%We look at the particular case $K = k(z)$ with the usual derivative $\frac{d}{dz}$ and algebraically closed $k$ of characteristic zero. The solution to the inverse problem over $k(z)$ has been known for some time, as \cite{inverse} gives a positive answer : every linear algebraic group over $k$ can indeed be realized as a differential Galois group over $k(z)$.
%
%\newpar
%One may also note that most inverse problems can actually be formulated as questions on differential Galois group of families of equations : for instance, the inverse problem over $k(z)$ can be reformulated as asking what the algebraic quotients of the absolute differential Galois group of $k(z)$ are, as the Galois groups of differential equations over $k(z)$ are exactly the algebraic quotients of the absolute differential Galois group.
%
%\newpar
%%We are interested in realizing algebraic groups as Galois groups of equations of a particular class : regular singular equations. When $k=\C$, it has been known for some time that every linear algebraic group over $\C$ can be realized as the Galois group of a regular singular differential equation. We can actually get an even better, more precise statement as a consequence of the Riemann-Hilbert correspondence : every linear algebraic group over $\C$ generated by $d$ elements can be realized as the Galois group of a regular singular equation with singularities in any $S \subseteq \P^1(\C)$ of cardinality $d+1$.
%
%\newpar
%\cite{fengwib} provides a new, short solution to the inverse problem for $k(z)$ over arbitrary $k$ using only the well-known solution over $\C$ and a specialization theorem. We claim that this proof can, with minor modifications, be extended to a solution for the regular singular inverse problem (Theorem \ref{thm:inv}) :
%
%\begin{theorem}
%	
%	Let $k$ be an algebraically closed field of characteristic zero, $G$ an algebraic group over $k$ generated by $d$ elements, and $S$ a subset of $\P^1(k)$ of cardinality $d+1$. Then $G$ is the Galois group of some regular singular matrix differential equation over $k(z)$ with singularities in $S$.
%	In particular, every linear algebraic group can be realized as the Galois group of a regular singular differential equation.
%\end{theorem}
%
%\noindent
%The general strategy is to notice that it is enough to prove the theorem for countable subfields of $\C$, and use the specialization theorem (Theorem \ref{thm:spe}) from \cite{fengwib} to take the result down from $\C$ to countable subfields. The main contribution we make is to show that if the singularities are in the subfield, a regular singular matrix equation remains regular singular under specializations.

%\section{A solution for the regular singular inverse problem}

\section{The regular singular inverse problem}

Throughout this article $k$ is an algebraically closed field of characteristic zero and all fields are assumed to have characteristic zero. All rings are assumed to be commutative. For the sake of clarity, we use $t$ for the local variable and $z$ for the global variable. In other words, we write $k((t))$ for the field of formal Laurent series over $k$ and $k(z)$ for the field of rational functions over $k$. For brevity, we use the term ``algebraic group (over $k$)'' for ``affine group scheme of finite type (over $k$)''. In other words, all algebraic groups are assumed to be linear.
A ``proalgebraic group (over $k$)'' is an affine group scheme (over $k$). A ``closed subgroup'' of a proalgebraic group is a closed subgroup scheme.
For a morphism $R\to R'$ of rings and a scheme $X$ over $R$ we denote with $X_{R'}$ the scheme over $R'$ obtained from $X$ by base change via $R\to R'$. For an affine scheme $X$ over a ring $R$, we denote with $R[X]$ the $R$-algebra representing $X$. In particular, for an algebraic group $G$ over $k$, its coordinate ring is denoted with $k[G]$.

\medskip

In this section we prove the results announced in the introduction. After recalling the basics of differential Galois theory and the theory of regular singular differential equations, we study regular singular differential equations under an extension of the constants and under specialization. These preparatory results are then combined to yield the proof of the main theorem.
	
\subsection{Differential Galois theory}
	
In this section we recall a few key definitions and facts from differential Galois theory. For more background see any of the textbooks \cite{SVDP}, \cite{Magdid:Lectures}, \cite{CrespoHajto} or \cite{Saul}.

\medskip

%
% More definitions, general results and facts can be found in (section 1 of) \cite{SVDP}. As a general convention, we refer to affine group schemes of finite type (over a field $k$) as algebraic groups (over $k$) - in particular, our algebraic groups are all linear. All fields will have characteristic zero.

Let $(K,\del)$ be a differential field with field of constants $K^\del=\{a\in K|\ \del(a)=0\}$ equal to $k$. A \emph{differential module} $(M,\nabla)$ over $(K,\del)$ is a finite dimensional $K$-vector space together with an additive map $\nabla\colon M\to M$ satisfying $\nabla(am)=\del(a)m+a\nabla(m)$ for $a\in K$ and $m\in M$.	
To a differential equation $\del(y)=Ay$ with $A\in K^{n\times n}$ we can associate the differential module $(K^n,\nabla_A)$, where $\nabla_A=\partial-A\colon K^n\to K^n$ is given by $\nabla_A(\xi)=\del(\xi)-A\xi$.

Conversely, if $(M,\nabla)$ is a differential module and $\underline{e}=(e_1,\ldots,e_n)$ is a $K$-basis of $M$, then
we can write $\nabla(\underline{e})=\underline{e}(-A)$ for some matrix $A\in K^{n\times n}$ so that 
$\nabla(\underline{e}\xi)=\underline{e}\del(\xi)+\nabla(\underline{e})\xi=\underline{e}(\del(\xi)-A\xi)$ for all $\xi\in K^n$. Thus, up to the choice of a basis, $(M,\nabla)$ is of the form $(K^n,\nabla_A)$. The choice of another basis $\underline{e}'$, say $\underline{e}=\underline{e}'F$ with $F\in\GL_n(K)$, leads to another matrix $A'\in K^{n\times n}$ satisfying
\begin{equation}\label{eq:gauge equivalence}
	A'=FAF^{-1}+\del(F)F^{-1} 		
\end{equation}	
Two matrices $A,A'\in K^{n\times n}$ are said to be \emph{Gauge equivalent} (over $K$) if there	exists a matrix $F\in\GL_n(K)$ such that (\ref{eq:gauge equivalence}) is satisfied.


Base change for differential modules works as expected: If $(K',\del')$ is a differential field extension of $(K,\del)$ and $(M,\nabla)$ is a differential module over $(K,\del)$, then $M'=M\otimes_{K}K'$ becomes as differential module over  $(K',\del')$ via $\nabla'(m\otimes a)=\nabla(m)\otimes a+m\otimes\del'(a)$ for $m\in M$ and $a\in K'$. If $(M,\nabla)=(K^n,\nabla_A)$ for some $A\in K^{n\times n}$, then $(M',\nabla')=(K'^n,\nabla_A)$. I.e., in terms of differential equations, this construction simply means that a differential equation $\del(y)=Ay$ over $K$ is considered as a differential equation over $K'$.


\begin{definition} \label{defi:PVring}
For $A \in K^{n \times n}$, a differential $K$-algebra $R$ is a \emph{Picard-Vessiot ring} for $\del(y) = Ay$ if
\begin{enumerate}[(i)]
	\item $R$ is $\del$-simple (i.e., it has no $\del$-stable ideals other than $\{0\}$ and $R$),
	\item there exists a fundamental solution matrix for $\del(y)=Ay$ in $R$, i.e., there exists $Y \in \GL_n(R)$ such that $\del(Y) = AY$,
		\item as a $K$-algebra $R$ is generated by the entries of $Y$ and the inverse of its determinant, i.e., $R = K\big[ Y_{i,j}, \frac{1}{\det(Y)}\big]$.                                                                                                                                         
\end{enumerate}
\end{definition}

Given $A\in K^{n\times n}$, there exists a unique (up to $K$-$\del$-isomorphisms) Picard-Vessiot ring $R/K$ for $\del(y)=Ay$. The \emph{differential Galois group} of $\del(y)=Ay$ or of $R/K$ is then defined as the group of differential automorphisms of $R/K$. More formally, we define $\Gal(R/K)$ as the functor, from the category of $k$-algebras to the category of groups, that associates to a $k$-algebra $T$ the group of differential automorphisms of $R\otimes_k T$ over $K\otimes_k T$, where $T$ is considered as a constant differential ring. The functor $\Gal(R/K)$ is representable by a finitely generated $k$-algebra, i.e., $\Gal(R/K)$ is an algebraic group (over $k$).



Definition \ref{defi:PVring} and the definition of the differential Galois group can be generalized to (potentially infinite) families $(\del(y)=A_iy)_{i\in I}$ of linear differential equations over $K$. For example, for $K=k(z)$ and $S$ a subset of $\P^1(k)$, one can consider the differential Galois group of the family of all regular singular differential equations over $k(z)$ with singularities contained in $S$ (see Section~\ref{subsec: Regular singular differential equations} below for precise definitions). In condition (ii) one requires that for every member of the family there exists a fundamental solution matrix in $R$ and in condition (iii) one allows all entries and the inverse of the determinant of all fundamental solution matrices. 

For families $\Gal(R/K)$ need not be an algebraic group, however, it is still a proalgebraic group (over $k$). Also note that the Picard-Vessiot ring for a family $(\del(y)=A_iy)_{i\in I}$ is generated by the Picard-Vessiot rings of its members $\del(y)=A_iy$. 

%\begin{definition}
%Let $(K,\del)$ be a differential field, $k$ its field of constants, $A \in K^{n\times n}$, and $R$ a Picard-Vessiot ring for the linear matrix differential equation $\del y=Ay$. The Galois group of (the Picard-Vessiot ring of) the equation is the group functor $\Gal(R/K) = \underline{\Aut}(R/K)$ defined on $k$-algebras as the group $\underline{\Aut}(R/K)(S) = \Aut^\del(R\otimes_k S/K\otimes_k S)$ of differential $K\otimes_k S$-algebra automorphisms of $R \otimes_k S$, where $S$ is taken as a constant $k$-algebra.
%\end{definition}

%\noindent
%The group functor $\Gal(R/K)$ is actually an algebraic group, and it is represented by the $k$-algebra $(R \otimes_K R)^\del$. If we take a fundamental matrix solution $Y$ to $\del y = Ay$ in $\GL_n(R)$, the algebra $(R \otimes_K R)^{\del}$ is generated over $k$ by the coefficients and the inverse of the determinant of $Y^{-1} \otimes Y$. We can again give an analogous definition for a family of differential equations, with the difference that $\Gal(R/K)$ will not necessarily be an algebraic group. Indeed, setting $C^\alpha := Y_\alpha^{-1} \otimes Y_\alpha$, $\Gal(R/K)$ is represented by $(R \otimes_K R)^\del = k\lb C_{i,j}^\alpha, \frac{1}{\det(C^\alpha)}\rb$, which is in general not finitely generated.
%	
%
%	
%\noindent
%We can give an analogous definition for a family of linear differential equations $(\del y = A_\alpha y)_{\alpha}$~: $R$ is a Picard-Vessiot ring for such a family of equation if it is a simple differential $K$-algebra that is generated over $K$ by the coefficients of fundamental matrix solutions to the equations $\del y = A_\alpha y$ and the inverses of the determinants. In that case, $R$ can be written as a union of smaller Picard-Vessiot rings, namely the Picard-Vessiot rings of finitely many of the equations.	

	
	\subsection{Regular singular differential equations} \label{subsec: Regular singular differential equations}
	
	
	In this section we recall the basic definitions regarding regular singular differential equations and observe that the relevant properties are preserved under an extension of the base field $k$. We also establish a key result (Lemma \ref{lem:basis}) for the proof of our main theorem (Theorem \ref{thm:inv}): We show that if a matrix $A\in k(t)^{n\times n}$ is Gauge equivalent over $k((t))$ to a matrix $A'$ such that $tA'\in k[[t]]^{n\times n}$, then $A$ is already Gauge equivalent over $k(t)$ to such a matrix. More background on regular singular differential equations can be found in \cite[Sections 5 and 6]{SVDP}, \cite{Saul} or \cite[Part~I]{MitschiSauzin}.
	
\medskip	
	
	We first recall the basic definitions regarding formal differential modules, i.e., differential modules over the differential field $(k((t)),\frac{d}{dt})$ of formal Laurent series with coefficients in $k$.
	
	
	
	
	 A \emph{$k[[t]]$-lattice} $\Lambda$ in a finite dimensional $k((t))$\=/vector space $M$ is a $k[[t]]$-submodule of the form $\Lambda=k[[t]]e_1+\ldots+k[[t]]e_n$, where $e_1,\ldots,e_n$ is a $k((t))$-basis of $M$.  
	A differential module $(M,\nabla)$ over $(k((t)),\frac{d}{dt})$ is 	
	\begin{itemize}
		 \item \emph{regular} if there exists a $k[[t]]$-lattice $\Lambda$ in $M$ such that $\nabla(\Lambda)\subseteq \Lambda$;
		 \item \emph{regular singular} if there exists a $k[[t]]$-lattice $\Lambda$ in $M$ such that $\nabla(\Lambda)\subseteq \frac{1}{t}\Lambda$.
	\end{itemize}
	
For a matrix $A\in k((t))^{n\times n}$, the differential equation $\del(y)=Ay$ is \emph{regular} (or \emph{regular singular}) if the corresponding differential module $(k((t))^n,\nabla_A)$ is regular (or regular singular). Thus $\del(y)=Ay$ is regular if and only if $A$ is Gauge equivalent over $k((t))$ to a matrix $A'\in k[[t]]^{n\times n}$ and $\del(y)=Ay$ is regular singular if and only if $A$ is Gauge equivalent over $k((t))$ to a matrix $A'\in k((t))^{n\times n}$ such that $tA'\in k[[t]]^{n\times n}$. 


\medskip

We next discuss our main topic of interest, linear differential equations over $k(z)$, the field of rational functions with coefficients in $k$. Throughout this article we consider $k(z)$ as a differential field with respect to the standard derivation $\frac{d}{dz}$. 

A differential module $(M,\nabla)$ over $k(z)$ defines for every $p\in\P^1(k)=k\cup\{\infty\}$ a formal differential module $(M_p,\nabla_p)$. In detail: For $p\in k$ we set $\widehat{k(z)}^p=k((z-p))$. Then, writing $t=z-p$ so that $\widehat{k(z)}^p=k((t))$, we see that $(k(z),\frac{d}{dz})$ is a differential subfield of $(k((t)),\frac{d}{dt})$. We define the differential module $(M_p,\nabla_p)$ over $\widehat{k(z)}^p=k((t))$ as the base change of $(M,\nabla)$ via the inclusion $k(z)\subseteq \widehat{k(z)}^p$. 

For $p=\infty$, we set $\widehat{k(z)}^p=k((z^{-1}))$. Then, writing $t=z^{-1}$ so that $\widehat{k(z)}^p=k((t))$, we see that $(k(z),\frac{d}{dz})$ is a differential subfield of $(k((t)),-t^2\frac{d}{dt})$. Thus, via base change from  $(k(z),\frac{d}{dz})$ to $(k((t)),-t^2\frac{d}{dt})$ we obtain from $(M, \nabla)$ a differential module $(M_p,\nabla')$ over  $(k((t)),-t^2\frac{d}{dt})$. Set $\nabla_p=-\frac{1}{t^2}\nabla'\colon M_p\to M_p$. Then $(M_p,\nabla_p)$ is a differential module over $(k((t)),\frac{d}{dt})$. 

A point $p\in\P^1(k)$ is \emph{regular} for $(M,\nabla)$ if $(M_p,\nabla_p)$ is regular. Otherwise it is \emph{singular}. A point $p\in\P^1(k)$ is \emph{regular singular} for $(M,\nabla)$ if $(M_p,\nabla_p)$ is regular singular. Finally, $(M,\nabla)$ is \emph{regular singular} if all points $p\in\P^1(k)$ are regular singular for $(M,\nabla)$. As regular points are regular singular, this is equivalent to saying that all singular points are regular singular. 

For $A\in k(z)^{n\times n}$, a point $p\in \P^1(k)$ is \emph{regular} (or \emph{regular singular}) for $\del(y)=Ay$ if $p$ is a regular (or regular singular) point of the corresponding differential module $(k(z)^n,\nabla_A)$. 
So, explicitly, a point $p\in k$ is regular (or regular singular) for $\del(y)=Ay$ if and only if $\del(y)=A(t+p)y$ is regular (or regular singular). On the other hand, $p=\infty$ is a regular (or regular singular) point for $\del(y)=Ay$ if and only if $\del(y)=-\frac{1}{t^2}A(\frac{1}{t})y$ is regular (or regular singular).



The \emph{singularities} or \emph{singular points} of $\del(y)=Ay$ are the points of $\P^1(k)$ that are not regular for $\del(y)=Ay$. A point $p\in k$ is a \emph{pole of $\del(y)=Ay$} if it is a pole of one of the entries of $A$. The point $p=\infty$ is \emph{a pole of $\del(y)=Ay$} if $t=0$ is a pole of $-\frac{1}{t^2}A(\frac{1}{t})$.

 If $p\in\P^1(k)$ is not a pole of $\del(y)=Ay$, then $A\in k[[t]]^{n\times n}$ for $t=z-p$ if $p\in k$ and $-\frac{1}{t^2}A(\frac{1}{t})\in k[[t]]^{n\times n}$ if $p=\infty$. Therefore, $p$ is a regular point for $\del(y)=Ay$. Thus, the singularities of $\del(y)=Ay$ are poles of $\del(y)=Ay$. On the other hand, a pole of $\del(y)=Ay$ need not be a singularity of $\del(y)=Ay$. For example, $p=0$ is a regular point for $\del(y)=\frac{1}{z}y$ (with solution $y(z)=z$). Indeed, for $F=\frac{1}{z}$, we have $F\frac{1}{z}F^{-1}+\del(F)F^{-1}=0$. In the literature, a pole of $\del(y)=Ay$ that is not a singularity of $\del(y)=Ay$ is sometimes referred to as an \emph{apparent singularity}. But note that, according to our terminology, an apparent singularity is not a singularity.
 
 
 The differential equation $\del(y)=Ay$ is \emph{Fuchsian} at $p\in\P^1(k)$ if the corresponding matrix has at most a pole of order one, i.e., $tA(t+p)\in k[[t]]^{n\times n}$ for $p\in k$ or $-\frac{1}{t}A(\frac{1}{t})\in k[[t]]^{n\times n}$ for $p=\infty$. If $\del(y)=Ay$ is Fuchsian at $p\in\P^1(k)$, then $\del(y)=Ay$ is regular singular at $p$. Conversely, if $\del(y)=Ay$ is regular singular at $p$, there exists a Gauge transformation over $k((t))$ leading to a matrix with at most a pole of order one. 
 
 




The differential equation $\del(y)=Ay$ is \emph{regular singular} if all points $p\in \P^1(k)$ are regular singular for $\del(y)=Ay$, i.e., if the differential module $(k(z)^n,\nabla_A)$ is regular singular. 

For $k=\mathbb{C}$, the condition of being regular singular is usually defined via a growth condition on the solutions (see e.g., \cite[Def. 9.11]{Saul} or \cite[Part I, Def. 1.24]{MitschiSauzin}). However, this growth condition is equivalent to the existence of a local meromorphic gauge transformation yielding a transformed matrix that has at most a pole of order one (\cite[Theorem 9.17]{Saul}). The existence of such a local meromorphic gauge transformation is equivalent to the existence of a local formal gauge transformation with the same effect (\cite[Prop. 3.12 and Lemma 3.42]{SVDP}). Thus the above definition of a regular singular differential equation $\del(y)=Ay$ with $A\in k(z)^{n\times n}$ is equivalent to the classical definition for $k=\mathbb{C}$. 


The following lemma shows that the conclusion of Theorem \ref{theo: main} remains unfazed under an extension of the base field $k$.
%
%\begin{lemma} \label{lemma: base change for formal}
%Let $k\subseteq k'$ be an inclusion of algebraically closed fields and let $(M,\nabla)$ be a differential module over $(k((t)), \frac{d}{dt})$. If $M$ is regular (or regular singular) then $M'=M\otimes_{k((t))}k'((t))$ is regular (or regular singular).
%\end{lemma}	
%\begin{proof}
%	If $\Lambda$ is a lattice in $M$ with $\nabla(\Lambda)\subseteq \Lambda$ (or $\nabla(\Lambda)\subseteq \frac{1}{t}\Lambda$), then $\Lambda'=k'[[t]]\Lambda$ is a lattice in $M'$ with $\nabla'(\Lambda')\subseteq \Lambda'$ (or $\nabla'(\Lambda')\subseteq \frac{1}{t}\Lambda'$). Thus if $M$ is regular (or regular singular), so is $M'$.
%\end{proof}		
%
%\begin{lemma} \label{lemma: base change for rational}
%	Let $k\subseteq k'$ be an inclusion of algebraically closed fields and let $(M,\nabla)$ be a differential module over $k(z)$. 
%	If $M$ is regular singular with singular points in $S\subseteq \P^1(k)$, then $M'=M\otimes_{k(z)}k'(z)$ is regular singular with singular points in $S$.
%\end{lemma}
%\begin{proof}
%	We may assume that $(M,\nabla)=(k(z)^n,\nabla_A)$ for some $A\in k(z)^{n\times n}$.
%Let $p\in\P^1(k')$ be a singular point for $M'$. We have to show that $p$ is regular singular for $M'$ and that $p\in S$. All points in $k'$ that are not poles of an entry of $A$ are regular for $M'$. Thus $p$ must be a pole of an entry of $A$ or $p=\infty$. In particular, $p\in\P^1(k)$. By assumption, $p$ is a regular singular point for $M$. So $p$ is a regular singular point for $M'$ by Lemma \ref{lemma: base change for formal}. Also, by the same lemma, $p$ cannot be a regular point for $M$ because then it would be a regular point for $M'$. Thus $p$ is a singular point for $M$ and therefore lies in $S$. 
%\end{proof}


\begin{lemma}\label{lem:basechange}
	Let $k \subseteq k'$ be an inclusion of algebraically closed fields and $S=\{p_1,\ldots,p_{d+1}\}$ a subset of $\P^1(k)$ with $d+1$ elements. If $A\in k(z)^{n\times n}$ is such that $\del(y)=Ay$ has all poles contained in $S$, 	
	is Fuchsian at $p_1,\ldots,p_d$, regular singular at $p_{d+1}$ and has differential Galois group $G$,
	then $\del(y)=Ay$, considered as a differential equation over $k'(z)$, has all poles contained in $S$, 	
	is Fuchsian at $p_1,\ldots,p_d$, regular singular at $p_{d+1}$ and has differential Galois group $G_{k'}$.
\end{lemma}
\begin{proof}
If $p\in\P^1(k')$ is a pole of $\del(y)=Ay$, then $p\in\P^1(k)$. So the poles of $\del(y)=Ay$ do not depend on whether we consider the equation over $k(z)$ or over $k'(z)$. If $\del(y)=Ay$ is Fuchsian (or regular singular) at $p\in \P^1(k)$, when considered as a differential equation over $k(x)$, then $\del(y)=Ay$ is also Fuchsian (or regular singular) at $p$, when considered as a differential equation over $k(x)$.

The statement about the differential Galois group follows from \cite[Lemma 3.2]{fengwib}.
\end{proof}	

	
%	
%\begin{definition}
%	Let $k$ be an algebraically closed field of characteristic zero. We consider $k(z)$ as a differential field with the usual derivation $\del = \frac{d}{dz}$. A matrix differential equation $\del y = Ay$ with $A \in k(z)^{n\times n}$ is said to be regular singular at $0$ if there exists a matrix $B \in k((z))^{n\times n}$ with at most a pole of order one at zero, and a matrix $F \in \GL_n\big(k((z))\big)$ such that $\del(F)F^{-1}+FAF^{-1} = B$. We say the matrix $A$ is regular singular (at $0$) if the associated matrix differential equation is.
%	If $p \in k$, we say that $A$ is regular singular at $p$ if $A(z-p)$ is regular singular at zero. We say that $A$ is regular singular at $\infty$ if $-\frac{1}{w^2}A(1/w)$ is regular singular at $w=0$.\\
%	We say that a differential equation $\del y = Ay$ is regular singular if it is regular singular at every $p \in \P^1(k) = k \cup \{\infty\}$.
%\end{definition}
%
%\noindent
%The condition $\del(F)F^{-1}+FAF^{-1} = B$ can be reformulated as $F\lp\del-A\rp F^{-1}y = \lp\del-B\rp y$ for all $y \in k((z))^n$, or just $F\lp\del-A\rp F^{-1}= \del-B$. We then say that $A$ is gauge equivalent to $B$, and $A \mapsto \del(F)F^{-1}+FAF^{-1}$ is a gauge transformation. Gauge transformations define a group action of $\GL_n\big(k((z))\big)$ (this can be seen either by direct computation or by using the differential operator formulation).
%
%\newpar
%We can formulate this condition in a more invariant way : let $M$ be a $k(z)$-module. A connection $\nabla : M \to \Omega^1_{k(z)/k} \otimes_{k(z)} M = Mdz$ is said to be regular singular at $p$ if the connection it induces on $M_p := M \otimes_{k(z)} k((z-p))$ is regular singular, i.e. there exists a $k[[z-p]]$-lattice $\L_p$ in $M_p$ with $\nabla(\L_p) \subseteq d(z-p) \otimes (z-p)^{-1} \L_p$. Choosing a coordinate on $\P^1(k)$ (either $z$ or $w = 1/z$) and a basis for $M$ allows us to write $\nabla = d - Adz$ with $A \in k(z)^{n \times n}$. The switch between $z$ and $w$ results in multiplying $A$ with a factor of $-w^2$ or $-z^2$ because $dw = -\frac{1}{z^2}dz$, and the basis change results in a gauge transformation of $A$.
%
%\newpar
%A matrix $A \in k(z)^{n\times n}$ is then said to be regular singular at $p$ if the connection $d - Adz$ is regular singular at $p$. In particular, at infinity, this translates into $d+\frac{1}{w^2}Adw$ being regular singular at $0$.
%

%\newpar
%We need to be careful when studying regular differential matrix equations : indeed, simply looking at the poles of a matrix isn't enough to tell whether or not a matrix is regular singular. For instance, the matrix
%
%$$A(z) = \bbm -1/z & -1/z^2 \\ 1 & 0\ebm$$
%is regular singular at zero. Indeed, it is the companion matrix of the regular singular differential operator $Ly = z^2\del^2 y+z\del y+y$. Rewriting it using the Euler derivative $\delta = z\del$, we have $Ly = \delta^2 y + y$, and the associated base change yields that $A$ is gauge equivalent to
%
%$$B(z) = \bbm 0 & -1/z \\ 1/z & 0\ebm$$
%which is clearly regular singular at $0$.
%\\



Our next goal is to show that if a matrix $A\in k(t)^{n\times n}$ is Gauge equivalent over $k((t))$ to a matrix $A'$ such that $tA'\in k[[t]]^{n\times n}$, then $A$ is already Gauge equivalent over $k(t)$ to such a matrix.
For the sake of clarity of the exposition, we separate the proof into a series of lemmas.

%The next few lemmas might not feel very new or deep to the reader, but we include proofs for completeness and to avoid too many references.

%\begin{lemma}\label{lem:val}
%	Let $K$ be a discretely valued field, $\l O_K$ its valuation ring and $\pi \in \l O_K$ a generator of the valuation ideal. Then if $\L \subseteq K^n$ is an $\l O_K$-lattice, i.e. of the form $\L = \l O_K e_1 + ... + \l O_K e_n$ where $(e_1,...,e_n)$ is a basis of $K^n$ and $e \in K^n$ is any vector, there exists an $a \in K^\times$ such that $ae \in \L$.
%\end{lemma}	
%\begin{proof}
%	Write $e = a_1e_1+...+a_ne_n$, and let $m := \min\{r \in \N : \forall i \le n, \pi^ra_i \in \l O_K\}$. This is well-defined because $K = \bigcup_{r \in \N} \pi^{-r} \l O_K$, and by definition we have $\pi^m e \in \L$.
%\end{proof}



\begin{lemma}\label{lem:val}
	If $\Lambda\subseteq k((t))^n$ is a $k[[t]]$-lattice and $E$ is a finite subset of $k((t))^n$, then there exists a non-zero $h\in k((t))$ such that $he\in\Lambda$ for all $e\in E$.
\end{lemma}	
\begin{proof}
	Let $e_1,\ldots,e_n$ be a $k((t))$-basis of $k((t))^n$ such that $\L=k[[t]]e_1+\ldots+k[[t]]e_n$. Fix $e\in E$ and write $e = h_1e_1+...+h_ne_n$ with $h_1,\ldots,h_n\in k((t))$. Let $m_e\in\mathbb{N}$ be such that $t^{m_e}h_i\in k[[t]]$ for $i=1,\ldots,n$.	Then $t^{m_e}e\in\L$. Let $m$ be the maximum of all $m_e$'s. Then $h=t^m$ has the required property.
\end{proof}


\begin{lemma} \label{lemma: scaling lattice}
Let $(M,\nabla)$ be a differential module over $(k((t)),\frac{d}{dt})$ and $\L\subseteq M$ a $k[[t]]$-lattice such that $\nabla(\L)\subseteq \frac{1}{t}\L$. Then $h\L\subseteq M$ is a $k[[t]]$-lattice with $\nabla(h\L)\subseteq \frac{1}{t}h\L$
for any non-zero $h\in k((t))$.
\end{lemma}	
\begin{proof}
 Clearly $h\Lambda$ is a $k[[t]]$-lattice in $M$. Furthermore, $\nabla(h \L) \subseteq \frac{d}{dt}(h)\L + h\nabla(\L) \subseteq \frac{d}{dt}(h)\L + \frac{1}{t}h\L $. So it suffices to show that $\frac{d}{dt}(h)\L\subseteq\frac{1}{t}h\L$. But $\frac{t\frac{d}{dt}(h)}{h}\in k[[t]]$ because the order of $t\frac{d}{dt}(h)$ is the order of $h$ unless $h\in k$. Thus $\frac{d}{dt}(h)\L=\frac{h}{t} \frac{t\frac{d}{dt}(h)}{h} \L \subseteq \frac{h}{t} \L=\frac{1}{t}h\L$.
\end{proof}	

As usual, we denote with $k[t]_{(t)}$ the localization of the ring $k[t]$ at the prime ideal $(t)$, i.e., $k[t]_{(t)}=\left\{\frac{h_1}{h_2}|\ h_1,h_2\in k[t],\ t \text{ does not divide } h_2\right\}$.



\begin{lemma} \label{lem:latt}
Let $e_1,\ldots,e_n$ be a $k(t)$-basis of $k(t)^n$. For $\L=k[[t]]e_1+\ldots+k[[t]]e_n\subseteq k((t))^n$ we have $\L\cap k(t)^n=k[t]_{(t)}e_1+\ldots+k[t]_{(t)}e_n$.
\end{lemma}	
\begin{proof}
	The inclusion $k[t]_{(t)}e_1+\ldots+k[t]_{(t)}e_n\subseteq \L\cap k(t)^n$ follows from $k[t]_{(t)}\subseteq k[[t]]$.
	For the reverse inclusion, consider $m\in\L\cap k(t)^n$. As $e_1,\ldots,e_n$ is a $k((t))$-basis of $k((t))^n$, we have $m=a_1e_1+\ldots+a_ne_n$ for uniquely determined $a_1,\ldots,a_n\in k((t))$. As $m\in \L$, it follows that $a_i\in k[[t]]$ and as $m\in k(t)^n$ is follows that $a_i\in k(t)$. Thus $a_i\in k[[t]]\cap k(t)=k[t]_{(t)}$ for $i=1,\ldots,n$ and so $m\in k[t]_{(t)}e_1+\ldots+k[t]_{(t)}e_n$.
\end{proof}	

%\begin{lemma}\label{lem:latt}
%	
%	Let $K \subseteq L$ be an inclusion of non-archimedean valued fields, $\l O_K$, $\l O_L$ the corresponding valuation rings, and assume that $K \cap \l O_L = \l O_K$. Let $M$ be a $K$-vector space of finite dimension $n \ge 1$.
%	Let $e_1,...,e_n$ be a basis of $M$, and $\L = \l O_L e_1 + ... + \l O_L e_n$ the $\l O_L$-lattice in $M \otimes_K L$ generated by this basis. Then $\L \cap M = \l O_K e_1 + ... + \l O_K e_n$, and is an $\l O_K$-lattice in $M$.
%	
%\end{lemma}
%
%\begin{proof}
%	
%	We identify $M$ with a submodule of $M \otimes_K L$ by $m \mapsto m \otimes 1$. Any $v \in \L \cap M$ can be uniquely written as a linear combination of $e_i$ with coefficients in $\l O_L$ (because it is in $\L$) and in $K$ (because it is in $M$). Since the $(e_i)_i$ form a $L$-basis of $M \otimes_K L$, both are equal and the coefficients are in $\l O_L \cap K = \l O_K$. Conversely, any element of $\l O_K e_1+...+\l O_K e_n$ is in $\L \cap M$.
%	
%\end{proof}

The following lemma is a variant of Nakayama's Lemma. For a proof see e.g., \cite[Chapter~X, Theorem 4.4]{Lang:Algebra}.

\begin{lemma}\label{lem:fingen}
	Let $R$ be a local ring with maximal ideal $\fk m$ and residue field $k=R/\fk m$. If $M$ is a free $R$-module of finite rank, then a tuple of elements from $M$ is an $R$-basis of $M$ if and only if their images are a $k$-basis of $M/\fk m M$.	\qed
\end{lemma}


%\begin{lemma}\label{lem:fingen}
%	Let $R$ be a local ring with maximal ideal $\fk m$, $M$ a free module of finite dimension $n$ over $R$ and $K := R/\fk m$. Then :
%	\begin{enumerate}[$(i)$]
%		\item The elements $(f_i)_{i \in I}$ generate $M$ over $R$ if and only if the images of the $f_i$ generate $M/\fk m M$ over $K$.
%		\item The elements $f_1,...,f_n \in M$ form a free basis of $M$ over $R$ if, and only if the images of the $f_i$ form a free basis of $M / \fk m M$ over $K$.
%	\end{enumerate}
%\end{lemma}
%
%\begin{proof}
%	$(i)$ : The direct implication is clear. Conversely, let $N$ be the submodule of $M$ generated by the $(f_i)_i$. Since every element of $M/\fk m M$ is in the image of $N$, we have $M = N + \fk m M$, which, by Nakayama's lemma, implies $M = N$, and the $f_i$ generate $M$.\\
%	$(ii)$ : Since $M/ \fk m M$ is a vector space of dimension $n$ over $K$, the direct implication is clear. Conversely, we already know by $(i)$ that $f_1,...,f_n$ generate $M$. Choose a free basis of $M$ to identify it with $R^n$ : this identifies $M/\fk m M$ with $K^n$. Then, $(f_1,...,f_n)$ is a free basis of $M$ if, and only if the matrix $F$ representing it is invertible, which is equivalent to asking $\det(F) \notin \fk m$. But since the images of the $f_i$ form a basis of $M/\fk mM$, the associated matrix $\bar{F} \in K^{n\times n}$ we must have $\det(\bar{F}) \ne 0$, meaning $\det(F) \in R - \fk m = R^\times$.
%\end{proof}



%We also need the so-called cyclic vector lemma, proved for instance in \cite{SVDP}, Proposition 2.9 :
%
%\begin{theorem}[Cyclic vector]\label{lem:cycl}
%	Let $K$ be a differential field with algebraically closed field of constants $k \ne K$. If $(M,D)$ is a differential $K$-module of finite dimension $n$, there exists a cyclic vector $e \in M$, i.e. such that $e,D(e),...,D^{n-1}(e)$ generate $M$ as a $K$-vector space.
%\end{theorem}


Recall that the condition of being regular singular is characterized via local formal gauge transformations. The following lemma shows that in fact it is enough to consider rational gauge transformations. We believe this lemma is well-known but we could not find a suitable reference. For statements of a similar spirit in the analytic context see \cite[Lemma 12.1]{Saul} and \cite[Lemma 3.42]{SVDP}.
%This next lemma establishes that if a regular singular matrix has rational functions as coefficients, then its regular singularity can be seen with gauge transformations in $k(z)$ instead of $k((z))$.



\begin{lemma}\label{lem:basis}
	Let $A \in k(t)^{n\times n}\subseteq k((t))^{n\times n}$ be such that $\del(y)=Ay$ is regular singular at $0$. Then there
	exists a matrix $F\in \GL_n(k(t))$ such that $A'=FAF^{-1}+\del(F)F^{-1}$ satisfies $tA'\in k[t]_{(t)}^{n\times n}$. In other words, if $A$ is Gauge equivalent over $k((t))$ to a matrix $A'$ satisfying $tA'\in  k[[t]]^{n\times n}$, then $A$ is already Gauge equivalent over $k(t)$ to such a matrix.
\end{lemma}
\begin{proof}
	As $\del(y)=Ay$ is regular singular at $0$, there exists a $k[[t]]$-lattice $\Lambda$ in $k((t))^n$ such that $\nabla_A(\Lambda)\subseteq \frac{1}{t}\Lambda$, where $\nabla_A=\frac{d}{dt}-A$. 
	%So $\nabla_\delta(\Lambda)\subseteq\Lambda$, where $\nabla_\delta= z\nabla_A=z \frac{d}{dz} - zA$. Note that $(k((z))^n,\nabla_\delta)$ is a differential module over the differential field $(k((z)), \delta)$, where $\delta= z \frac{d}{dz}$ is the Euler derivative.
	
%%	By the cyclic vector lemma (\cite[Prop. 2.9]{SVDP}), there exists a cyclic vector for the differential module $(k(z)^n,\nabla_\delta)$ over the differential field $(k(z),\delta)$, i.e., there exists $e\in k(z)^n$ such that $e,\nabla_\delta(e),\nabla_\delta^2(e),\ldots,\nabla_\delta^{n-1}(e)$ is a $k(z)$-basis of $k(z)^n$. 
%	As a $k(z)$-basis of $k(z)^n$ is also a $k((z))$\=/basis of $k((z))^n$,
%	%(e.g., because this can be tested by the non-vanishing of an $n\times n$-determinant),
%	 we see that $e$ is also a cyclic vector for the differential module $(k((z))^n,\nabla_\delta)$ over the differential field $(k((z)),\delta)$.
%	

Let $\underline{e}=(e_1,\ldots,e_n)$ be the standard basis of $k(t)^n$. By Lemma \ref{lem:val} there exists a non-zero $h\in k((t))$ such that $he_i\in \L$ for $i=1,\ldots,n$. Then $\frac{1}{h}\L$ is a $k[[t]]$-lattice in $k((t)))^n$ with $\nabla_A(\frac{1}{h}\L)\subseteq \frac{1}{t}\frac{1}{h}\L$ (Lemma~\ref{lemma: scaling lattice}) and $e_1,\ldots,e_n\in \frac{1}{h}\L$.
Replacing $\L$ with $\frac{1}{h}\L$, we can thus assume that $e_1,\ldots,e_n\in \L$.
	
	Let $\L'$ be the  $k[[t]]$-submodule of $\L$ generated by all $(t\nabla_A)^i(e_j)$, where $i\geq 0$, and $j\in \{1,\ldots,n\}$. Note that, by construction, $t\nabla_A(\L')\subseteq \L'$.
	
	
	
	Recall (e.g. from \cite[Chapter III, Theorem 7.1]{Lang:Algebra}) that a submodule of a free module over a principal ideal domain is free of rank bounded by the rank of the ambient module. As $k[[t]]$ is a principal ideal domain and $\Lambda$ is a free $k[[t]]$-module of rank $n$, it follows that $\L'$ is a free $k[[t]]$-module of rank at most $n$.
	As $\L'$ contains the standard basis $\underline{e}$, we see that in fact $\L'$ is free of rank $n$.
	So $\L'/t\L'$ is a $k$-vector space of dimension $n$.
	% Since the $(t\nabla_A)^i(e_j)$'s generate $\L'$ as a $k[[t]]$-module, their images generate  $\L'/t\L'$ as a $k$-vector space. We can select $n$ of these generators, say $e_1',\ldots,e_n'$, such that their images are $k$-basis of $\L'/t\L'$. It follows from Lemma~\ref{lem:fingen}  that $e_1',\ldots,e_n'$ are a $k[[t]]$-basis of $L'$. In particular,  $e_1',\ldots,e_n'$ are $k[[t]]$-linearly independent and they lie in $k(t)^n$. Thus $e_1',\ldots,e_n'$ are a $k(t)$-basis of $k(t)^n$. 
	
	As the $(t\nabla_A)^i(e_j)$'s generate $\L'$ as a $k[[t]]$-module, Lemma \ref{lem:fingen} shows that we can select $n$ of these generators, say $e_1',\ldots,e_n'$, such that they form a $k[[t]]$-basis of $\L'$. In particular,  $e_1',\ldots,e_n'$ are $k[[t]]$-linearly independent and they lie in $k(t)^n$. Thus $\underline{e}=(e_1',\ldots,e_n')$ is a $k(t)$-basis of $k(t)^n$. 

	Lemma \ref{lem:latt} applied to $e_1',\ldots,e_n'$ and $\L'=k[[t]]e_1'+\ldots+k[[t]]e_n'$ shows that 
	$\L'\cap k(t)^n=k[t]_{(t)}e'_1+\ldots+k[t]_{(t)}e'_n$. As $\L'$ and $k(t)^n$ are stable under $t\nabla_A$, it follows that also $k[t]_{(t)}e'_1+\ldots+k[t]_{(t)}e'_n$ is stable under $t\nabla_A$. That is, if we write $\nabla_A(\underline{e}')=\underline{e'}(-A')$, then $tA'\in k[t]_{(t)}^{n\times n}$. For $F\in\GL_n(k(t))$ with $\underline{e}=\underline{e}'F$, we have $A'=FAF^{-1}+\del(F)F^{-1}$ as desired.
	%	We claim that $f\Lambda$ is a $\nabla_\delta$-invariant lattice for any non-zero $f\in k((z))$. Clearly $f\Lambda$ is a lattice in $k((z))^n$. Furthermore, $\nabla_\delta(f \L) \subseteq \delta(f)\L + f\nabla_\delta(\L) \subseteq \delta(f) \L + f\L $. So it suffices to show that $\delta(f) \L\subseteq f\L$. But $\frac{\delta(f)}{f}\in k[[z]]$ because that order of $\delta(f)$ is the order of $f$ unless $f\in k$. Thus $\delta(f)\L=f\frac{\delta(f)}{f} \L \subseteq f \L$.
%%	
%	
%
%	
%	
%	
%	\vspace{30mm}
%	
%
%%	Then, for any $f \in k((z))^\times, f\L$ is again a $\nabla_\delta$-invariant lattice. Indeed, $\nabla_\delta(f \L) \subseteq \delta(f)\L + f\nabla_\delta(\L) \subseteq \delta(f) \L + f\L $.\\
%%	Notice that $v(\delta(f)) \ge v(f)$ where $v$ is the $z$-adic valuation on $k((z))$, meaning that $v(\delta(f)/ f) = v(\delta(f)) - v(f) \ge 0$, i.e. $\delta(f) / f \in k[[z]]$. This finally means that $\delta(f)\L = f \frac{\delta(f)}{f} \L \subseteq f \L$.\\
%%	A cyclic vector $e$ exists for $k(z)^n$ with the derivative $\nabla_\delta$ (by lemma \ref{lem:cycl})- such a vector is also cyclic for $k((z))^n$ for the same derivative.
%%	Using lemma \ref{lem:val}, we can assume that $e \in \L$ by replacing $\L$ with some $f \L$, which is also $\nabla_\delta$-stable.\\
%	Since $e$ is a cyclic vector, $k((z))^n$ is spanned by $\{\nabla_\delta^m e : m \in \N\}$. The $k[[z]]$-submodule $\L'$ of $\L$ generated by the $\nabla_\delta^m e$ generates $k((z))^n$ as a $k((z))$-vector space, and it is free of rank $\le n$ because $k[[z]]$ is a principal ideal domain. Since $\L'$ spans $k((z))^n$ as a vector space, we deduce that it must be of rank $n$ : it is therefore a $k[[z]]$-lattice, and it is $\nabla_\delta$-stable.\\
%	Using lemma \ref{lem:fingen}, we deduce that there exists $i_1,...,i_n$ such that $\nabla_\delta^{i_1} e,..., \nabla_\delta^{i_n} e$ are a free basis of $\L'$ over $k[[z]]$. Because $\nabla_\delta^j e \in k(z)^n$ for all $j\ge0$, we may now use lemma \ref{lem:latt} with $K = k(z)$ embedded into its completion $L = k((z))$ with respect to the $z$-adic absolute value. We know that $k(z) \cap k[[z]]$ is $\l O$, and thus $\L' \cap k(z)^n$ is a $\nabla_\delta$-stable $\l O$-lattice in $k(z)^n$. The matrix of $\nabla_\delta$ in a basis of such a lattice is thus a regular matrix, and it is equivalent to $A$ by construction.
\end{proof}	






%\begin{lemma}\label{lem:basis}
%	?? Need this also for regular?/
%	Let $k$ be an algebraically closed field of characteristic zero, $A \in k(z)^{n\times n}$ be a matrix that is regular singular at $0$. Then, there exists a basis of $k(z)^n$ in which $\frac{d}{dz} - A$ is of the form $\frac{d}{dz} - B$ where $zB$ is a regular matrix, i.e. it has coefficients in $\l O := k[z]_{(z)}$. Equivalently, there exists a matrix $F \in \GL_n\big(k(z)\big)$ such that $F\lp\frac{d}{dz} - A\rp F^{-1} = \frac{d}{dz} - B$.
%\end{lemma}
%
%
%\begin{proof}
%	Saying that $A$ is regular singular (at $0$) implies the existence of a $k[[z]]$-lattice $\Lambda$ in $k((z))^n$ that is invariant with respect to the derivation $\nabla_\delta = z \frac{d}{dz} - zA$.\\
%	Notice that this derivative makes $k((z))^n$ into a differential $k((z))$-module for the Euler derivative $\delta = z \frac{d}{dz}$.
%	Then, for any $f \in k((z))^\times, f\L$ is again a $\nabla_\delta$-invariant lattice. Indeed, $\nabla_\delta(f \L) \subseteq \delta(f)\L + f\nabla_\delta(\L) \subseteq \delta(f) \L + f\L $.\\
%	Notice that $v(\delta(f)) \ge v(f)$ where $v$ is the $z$-adic valuation on $k((z))$, meaning that $v(\delta(f)/ f) = v(\delta(f)) - v(f) \ge 0$, i.e. $\delta(f) / f \in k[[z]]$. This finally means that $\delta(f)\L = f \frac{\delta(f)}{f} \L \subseteq f \L$.\\
%	A cyclic vector $e$ exists for $k(z)^n$ with the derivative $\nabla_\delta$ (by lemma \ref{lem:cycl})- such a vector is also cyclic for $k((z))^n$ for the same derivative.
%	Using lemma \ref{lem:val}, we can assume that $e \in \L$ by replacing $\L$ with some $f \L$, which is also $\nabla_\delta$-stable.\\
%	Since $e$ is a cyclic vector, $k((z))^n$ is spanned by $\{\nabla_\delta^m e : m \in \N\}$. The $k[[z]]$-submodule $\L'$ of $\L$ generated by the $\nabla_\delta^m e$ generates $k((z))^n$ as a $k((z))$-vector space, and it is free of rank $\le n$ because $k[[z]]$ is a principal ideal domain. Since $\L'$ spans $k((z))^n$ as a vector space, we deduce that it must be of rank $n$ : it is therefore a $k[[z]]$-lattice, and it is $\nabla_\delta$-stable.\\
%	Using lemma \ref{lem:fingen}, we deduce that there exists $i_1,...,i_n$ such that $\nabla_\delta^{i_1} e,..., \nabla_\delta^{i_n} e$ are a free basis of $\L'$ over $k[[z]]$. Because $\nabla_\delta^j e \in k(z)^n$ for all $j\ge0$, we may now use lemma \ref{lem:latt} with $K = k(z)$ embedded into its completion $L = k((z))$ with respect to the $z$-adic absolute value. We know that $k(z) \cap k[[z]]$ is $\l O$, and thus $\L' \cap k(z)^n$ is a $\nabla_\delta$-stable $\l O$-lattice in $k(z)^n$. The matrix of $\nabla_\delta$ in a basis of such a lattice is thus a regular matrix, and it is equivalent to $A$ by construction.
%\end{proof}	





\subsection{Generating algebraic groups}

In this short section we consider algebraic groups generated by finitely many elements and how this property behaves under base change.

Let $G$ be an algebraic group over $k$ and $g_1,\ldots,g_d\in G(k)$. The intersection of all closed subgroups $H$ of $G$ such that $g_1,\ldots,g_d\in H(k)$, is a closed subgroup of $G$ that also satisfies this property. Thus there exists a smallest closed subgroup of $G$ whose $k$-points contain $g_1,\ldots,g_d$. If this closed subgroup agrees with $G$, we say that $g_1,\ldots,g_d$ \emph{generated $G$ (as an algebraic group)}. 

Let $\Gamma$ be the (abstract) subgroup of $G(k)$ generated by $g_1,\ldots,g_d$. As the Zariski closure of $\Gamma$ is a subgroup of $G(k)$ (\cite[Lemma 2.2.4]{Springer:LinearAlgebraicGroups}), we see that $g_1,\ldots,g_d$ generate $G$ if and only if $\Gamma$ is Zariski dense in $G$. This means that $g_1,\ldots,g_d$ generate $G$ if and only if $f(\gamma)=0$ for all $\gamma\in\Gamma$ implies $f=0$ for $f\in k[G]$.



\begin{lemma}\label{lem:groupgen} Let $k \subseteq k'$ be an inclusion of algebraically closed fields, $G$ an algebraic group over $k$ and $g_1,\ldots,g_d\in G(k)$. Then 
	$g_1,\ldots,g_d$ generate $G$ if and only if $g_1,\ldots,g_d$ generate $G_{k'}$.
\end{lemma}
\begin{proof}
	Let $\Gamma$ be the (abstract) subgroup of $G(k)\subseteq G(k')=G_{k'}(k')$ generated by $g_1,\ldots,g_d$.
	First assume that $g_1,\ldots,g_d$ generate $G$. We have to show that $f'(\gamma)=0$ for all $\gamma\in\Gamma$ implies $f'=0$ for $f'\in k'[G_{k'}]=k[G]\otimes_k k'$. Let $(\lambda_i)_{i\in I}$ be $k$-basis of $k'$. Then $f'$ can be written as $f'=\sum f_i\otimes\lambda_i$ with $f_i\in k[G]$. For $\gamma\in \Gamma$, we have $0=f'(\gamma)=\sum f_i(\gamma)\lambda_i$. As the $\lambda_i$ are $k$-linearly independent and $f_i(\gamma_i)\in k$, it follows that $f_i(\gamma)=0$ for all $\gamma\in\Gamma$ and $i\in I$. Because $g_1,\ldots,g_d$ generate $G$, this implies that $f_i=0$. So $f'=0$ as desired.
	%	Write $\m I(G)$ the Hopf ideal of $G$, one has $\m I(G_{k'}) = \m I(G) \otimes_k k'$.\\
	%	For the direct implication, we want to prove that if a polynomial $\in k'[\GL_n] = k[\GL_n] \otimes_k k'$ is zero on all elements $\gamma \in \Gamma$, then it belongs to $\m I(G_{k'})$.Let $Q$ be such a polynomial, write $Q = \sum_{i=1}^r Q_i \otimes \lambda_i$ with $\lambda_i$ free over $k$ and $Q_i \in k[\GL_n]$. Then for all $\gamma \in \Gamma$, we have $0 = Q(\gamma) = \sum Q_i(\gamma) \lambda_i$, meaning $Q_i(\gamma) = 0$ for all $i$. Since $\Gamma$ is Zariski-dense in $G$, the $Q_i$ are all in $\m I(G)$ and we obtain $Q \in \m I(G)\otimes_k k' = \m I(G_{k'})$, proving that $\Gamma$ is Zariski-dense in $G_{k'}$.\\
	
	Conversely, assume that $g_1,\ldots,g_d$ generate $G_{k'}$. Then an $f'\in k'[G_{k'}]=k[G]\otimes_k k'$ that vanishes on all of $\Gamma$ must be zero. In particular, this holds for $f\in k[G]$ and therefore $g_1,\ldots,g_d$ generate $G$.
	%
	%
	%
	%	The converse is obtained by noticing that $G(k)$ can be identified with a subset of $G_{k'}(k') = G(k')$ in a way that is compatible with polynomial evaluation : if we consider elements of $\Gamma$, $G(k)$ and $G(k')$ as actual matrices, then all we are saying is that a polynomial $P \in k[\GL_n] \subseteq k'[\GL_n]$ that vanishes on $\Gamma$ vanishes on all of $G(k')$, so it vanishes in particular on all of $G(k)$.
\end{proof}
%
%A consequence of the above lemma is that the minimal number of elements needed to generate an algebraic group can only go down under base change.








\subsection{Specializations}

The proof strategy for our main result is to deduce it from the already known special case $k=\mathbb{C}$ via a specialization argument. In this section we recall the specialization theorem from \cite{fengwib} required for this approach. We also show that the conclusion of Theorem \ref{theo: main} is preserved under specialization in an appropriate sense.


%The core theorem we need to exploit the already known result over $\C$ is the specialization theorem from \cite{fengwib}, and we need a few more definitions in order to state it : notably, we recall the definition of differential group actions and differential torsors.

We begin by recalling some terminology from \cite{fengwib}. Let $\l B$ be a $k$-algebra (considered as a constant differential ring). The specializations we will be considering are morphisms $c\in \Hom_k(\l B,k)$, i.e., $c\colon \l B\to k$ is a morphism of $k$-algebras. We consider $\l B[z]$ as a differential ring with respect to the derivation $\frac{d}{dz}$. For a monic polynomial $f\in \l B[z]$ a specialization $c\colon \l B\to k$ extends to a morphism $c\colon \l B[z]_f\to k(z)$ of differential $k$-algebras via $c(z)=z$.
For a matrix $\l A\in \l B[z]_f^{n\times n}$ we denote with $A^c\in k(z)^{n\times n}$ the matrix obtained from $\l A$ by applying $c\colon \l B[z]_f\to k(z)$ to the entries of $\l A$. Similarly, for $a\in \l B$ we will write $a^c$ rather than $c(a)$.


 Assuming that $\l B$ is an integral domain, so that we can consider the algebraic closure $K$ of the field of fractions of $\l B$, the specialization problem is about comparing the differential Galois group of the generic equation $\del(y)=\l A y$ (over $K(z)$) with the differential Galois group of the specialized equations $\del(y)=A^cy$ (over $k(z)$). To capture the idea of specializing differential Galois groups, we need to consider group schemes over $\l B$ and we also need some more terminology.





 Let $\l B$ be a $k$-algebra and let $\l Q$ be a differential $\l B$-algebra. (The most relevant case for us is $\l Q=\l B[z]_f$.) For a differential $\l Q$-algebra $\l R$, we define $\underline{\Aut}(\l R/\l Q)$ as the functor form the category of $\l B$-algebras to the category of groups given by $\underline{\Aut}(\l R/\l Q)(\l T) = \Aut^\del(\l R \otimes_{\l B} \l T / \l Q \otimes_{\l B} \l T)$, where $\l T$ is considered as a constant differential ring. On morphisms, $\underline{\Aut}(\l R / \l Q)$ is defined by base change.
An \emph{action} of an affine group scheme $\l G$ over $\l B$ on $\l R/\l Q$ is a morphism of group functors $\l G \to \underline{\Aut}(\l R/\l Q)$.

%\begin{definition}[Algebraic group actions]
%	
%	Let $k$ be an algebraically closed field of characteristic zero, $\l B$ a $k$-algebra, and $\l Q$ a differential $\l B$-algebra.
%	
%	\noindent
%	For a differential $\l Q$-algebra $\l R$, we define $\underline{\Aut}(\l R/\l Q)$ as the functor on $\l B$-algebras given by $\underline{\Aut}(\l R/\l Q)(\l T) := \Aut^\del(\l R \otimes \l T / \l Q \otimes \l T)$ where $\l T$ is taken as a constant ring. On morphisms, $\underline{\Aut}(\l R / \l Q)$ is defined by base change.
%	
%	\noindent
%	A (differential) action of an affine group scheme $\l G$ over $\l B$ on $\l R/\l Q$ is a morphism of group functors $\l G \to \underline{\Aut}(\l R/\l Q)$.
%	
%\end{definition}

Let $\l B[\l G]$ denote the coordinate ring of $\l G$, i.e., the $\l B$-algebra representing $\l G$. By \cite[Lemma 3.3]{fengwib}, to specify an action of $\l G$ on $\l R / \l Q$ is equivalent to specifying a morphism (the coaction) $\rho\colon \l R \to \l R \otimes_{\l B} \l B[\l G]$ of $\l Q$-$\del$-algebras such that 
\[\begin{tikzcd} \l R \otimes_{\l B} \l B[\l G] \otimes_{\l B}\l B[\l G] & \l R \otimes_{\l B} \l B[\l G] \arrow[swap]{l}{\rho \otimes \text{id}}\\ \l R \otimes_{\l B} \l B[\l G] \arrow{u}{\text{id}\otimes \Delta} & \l R \arrow[swap]{u}{\rho} \arrow{l}{\rho} \end{tikzcd}\qquad \text{ and }\qquad \begin{tikzcd}  \l R \otimes_{\l B} \l B[\l G] \arrow[swap]{d}{\text{id} \otimes \varepsilon} & \l R \arrow[swap]{l}{\rho} \\ \l R \otimes_{\l B} \l B \arrow[ur, equal]\end{tikzcd}\]
 commute, where $\Delta$ and $\varepsilon$ denote the comultiplicaiton and counit of the Hopf-algebra $\l B[\l G]$.





\begin{definition} \label{def: differential torsor}
%	Let $k$ be an algebraically closed field of characteristic zero, $\l B$ a $k$-algebra, $\l Q$ a differential $\l B$-algebra, $\l R$ a differential $\l Q$-algebra, $\l G$ an affine group scheme over $\l B$ and assume $\l G$ acts on $\l R / \l Q$.
%	\\
Assume, as above, that $\l G$ acts on $\l R / \l Q$. Then $\l R / \l Q$ is a \emph{differential $\l G$-torsor} if the morphism $\l R \otimes_{\l Q} \l R \to \l R \otimes_{\l B} \l B[\l G]$, $\ a \otimes b \mapsto (a \otimes 1) \cdot \rho(b)$ is an isomorphism. Moreover, for $\l A \in \l Q^{n \times n}$, if there exists a matrix $\l Y \in \GL_n(\l R)$ such that
	\begin{itemize}
		\item $\del(\l Y) = \l A \l Y$,
		\item $\l R = \l Q\lb \l Y_{i,j}, \frac{1}{\det(\l Y)}\rb$, and
		\item for every $\l B$-algebra $\l T$ and every $g \in \l G(\l T)$, there exists a matrix $M_g \in \GL_n(\l T)$ such that $g(\l Y \otimes 1) = (\l Y \otimes 1)(1 \otimes M_g)$,
	\end{itemize}
	 then $\l R / \l Q$ is a \emph{differential $\l G$-torsor for $\del(y) = \l Ay$}.
\end{definition}

The definition of a differential torsor can be interpreted geometrically: With $\l Z = \Spec(\l R)$, the coaction $\rho\colon \l R \to \l R \otimes_{\l B} \l B[\l G]=\l R\otimes_{\l Q}\l Q\otimes_{\l B }\l B[\l G]$ induces a group action $\l Z \times_{\l Q} \l G_{\l Q} \to \l Z$ (in the category of $\l Q$-schemes) and $\l R/\l Q$ is a differential $\l G$-torsor if and only if $\l Z \times_{\l Q} \l G_{\l Q} \to \l Z \times_{\l Q} \l Z$ is an isomorphism.
% of schemes over $\l Q$. In particular, for any $\l Q$-algebra $\l S$, $\l Z(\l S)$ is a $\l G_{\l Q}(\l S)$-torsor in the usual group-theoretic sense.

Definition \ref{def: differential torsor} mimics and generalizes the definition of a Picard-Vessiot ring.
In particular, if $R/K$ is a Picard-Vessiot ring for $\del(y)=Ay$ (Definition \ref{defi:PVring}), then $R/K$ is a 
differential $\Gal(R/K)$-torsor for $\del(y) = A y$.


The following specialization theorem is the key to our specialization argument.

\begin{theorem}[\cite{fengwib}, Theorem 4.26] \label{thm:spe}
	Assume that
	\begin{itemize}
		%\ibul $k$ is an algebraically closed field of characteristic zero,
		\ibul $\l B$ is a finitely generated $k$-algebra and an integral domain,
		\ibul $K$ is the algebraic closure of the field of fractions of $\l B$,
		\ibul $\l G$ is an affine group scheme of finite type over $\l B$,
		\ibul $f \in \l B[z]$ is a monic polynomial,
		\ibul $\l A \in \l B[z]^{n \times n}_f$,
		\ibul $\l R/\l B[z]_f$ is a differential $\l G$-torsor for $\del(y) = \l A y$ such that $\l R$ is flat over $\l B[z]_f$,
		\ibul $\l R \otimes_{\l B[z]_f} K(z)$ is $\del$-simple.
	\end{itemize}
	% $R^{\textup{gen}} =
	%\noindent
	Then there exists a $c \in \Hom_{k}(\l B, k)$ such that $R^c=\l R \otimes_{\l B[z]_f} k(z)$ is a Picard-Vessiot ring for $\del(y)=A^cy$ with differential Galois group $\l G_k$, where the involved base change is $c\colon \l B \to k$.
	
\end{theorem}

Note that the assumptions in Theorem \ref{thm:spe} imply that $R^{\textup{gen}}=\l R \otimes_{\l B[z]_f} K(z)$ is a Picard-Vessiot ring for $\del(y)=\l A y$ (over $K(z)$) with differential Galois group $\l G_K$. The following lemma shows that, starting from an inclusion $k\subseteq k'$ of algebraically closed fields and a differential equation $\del(y)=Ay$ over $k'(z)$, the technical conditions of Theorem \ref{thm:spe} can always be achieved.


%
%\newpar
%We introduce a final lemma from \cite{fengwib} that will let us spread out Picard-Vessiot rings, i.e. write them as extensions of scalars of differential torsors over finitely generated algebras. This allows for the use of the specialization theorem to specialize solutions of differential equations from bigger fields to their subfields.



\begin{lemma}[\cite{fengwib}, Lemma 3.11] \label{lem:spread}
	Let $k \subseteq k'$ be an inclusion of algebraically closed fields and let $A\in k'(z)^{n\times n}$. 
	%If $R/k'(z)$ is a Picard-Vessiot ring for $\partial(y) = A'y$ with fundamental solution matrix $Y' \in \GL_n(R')$,
	 Then there exist
	
	\begin{itemize}
		\item [$\bullet$] a finitely generated $k$-subalgebra $\l B$ of $k'$,
		\item[$\bullet$] a monic polynomial $f \in \l B[z]$,
		\item[$\bullet$] an affine group scheme $\l G$ of finite type over $\l B$, and
		\item[$\bullet$] a differential $\l G$-torsor $\l R/\l B[z]_f$,
	\end{itemize}
	such that
	\begin{itemize}
		\ibul $A  \in \l B[z]_f^{n\times n}$,
		\ibul $\l R$ is flat over $\l B[z]_f$ and $\l R / \l B[z]_f$ is a differential $\l G$-torsor for $\del(y) = Ay$
		\ibul $\l R \otimes_{\l B[z]_f} K(z)$ is $\del$-simple, where $K$ the algebraic closure of the field of fractions of $\l B$.
	%	\ibul $R^{\textup{gen}}\otimes_{K(z)} k'(z) \simeq R'$ by identifying $\l Y$ with $Y'$.
	\end{itemize}
	Moreover, for a fixed finitely generated $k$-algebra $\l B_0$ and monic $f_0 \in \l B_0[z]$, $\l B$ and $f$ can be chosen such that $\l B_0 \subseteq \l B$ and $f_0 $ divides $f$.
	Furthermore, if there exists an algebraic group $G$ over $k$ such that the differential Galois group of $\del(y)=Ay$ (over $k'(z)$) is of the form $G_{k'}$, then $\l G$ can be chosen to be $G_{\l B}$.
\end{lemma}



%
%\newpar
%The idea of the theorem is that if $R^{\text{gen}}/K(z)$ is a Picard-Vessiot ring for the "generic" differential equation $\del y = \l A y$, then for most specializations $c : \l B \to k$, the specialized ring $R^c$ will be a Picard-Vessiot ring for the specialized differential equation $\del y = A^c y$, and the Galois group of $R^c/k(z)$ is the "specialized" version of the Galois group of $R^{\text{gen}}/K(z)$.
%
%
%\newpar
%Using this, \cite{fengwib} shows that any algebraic group $G$ over a subfield $k$ of $\C$ such that $G_\C$ can be realized as a Galois group of a differential equation over $\C(z)$ (that is to say, all of them) can be realized as a Galois group of a differential equation over $k(z)$. Our main point is that this process preserves the regular singular nature of the equations and the amount of generators.
%
%\newpar
%The main issue we encounter in this endeavor is that the gauge matrices realizing the regular singularity of the equation have coefficients in $\C((z))$. Since the specialization theorem is true over finitely generated algebras, we might not be able to specialize every coefficient of the matrices at once. We will show that considering only gauge transformations with coefficients in $\C(z)$ is enough to establish the regular nature of the singularities of a matrix in $\C(z)^{n\times n}$, reducing the amount of coefficients we want to specialize to a finite one.
%


%Recall that $t$ is always the local variable (at zero).??


We next study regular singular differential equations under specialization, first locally, then globally.

\begin{lemma}\label{lem: specialize regsing at 0}
	Let $k\subseteq k'$ be an inclusion of algebraically closed fields and let $A \in k'(t)^{n \times n}\subseteq k'((t))^{n\times n}$ be such that $\del(y)=Ay$ is regular singular at $0$. Then there exists a finitely generated $k$\=/subalgebra $\l B$ of $k'$ and a monic polynomial $f \in \l B[t]$ such that $A \in \l B[t]_f^{n \times n}$ and $\del(y)=A^c y$ is regular singular at $0$ for every $c \in \Hom_k(\l B, k)$.	
\end{lemma}
\begin{proof}
By Lemma \ref{lem:basis}, there exists a matrix $F\in\GL_n(k'(t))$ such that the matrix
		$t(FAF^{-1}+\del(F)F^{-1})$ has entries in $k'[t]_{(t)}$. We may choose a monic polynomial $f\in k'[t]$ such that $fA$ and $fF$ have entries in $k'[t]$. Let $\l B$ be the $k$\=/subalgebra of $k'$ generated by the coefficients of $f$ and all coefficients of all entries of $fA$ and $fF$. Then $A,F\in \l B[t]_f^{n\times n}$.
		As $\det(F)\in \l B[t]_f$ is non-zero, we may write $\det(f)=\frac{h}{f^m}$ for some non-zero $h\in \l B[t]$. Adding the inverse of the leading coefficient of $h$ to $\l B$ and replacing $f$ with the product of $f$ with $h$ and the inverse of the leading coefficient of $h$, we may assume that $A, F, F^{-1}\in \l B[t]_f^{n\times n}$. 
		
		As $t(FAF^{-1}+\del(F)F^{-1})$ has entries in $k'[t]_{(t)}$, there exist non-zero elements $a_1,
		\ldots,a_m\in k'$ and a matrix $B\in k'[t]$ such that 
		\begin{equation} \label{eq: for regular}
					(t-a_1)\ldots(t-a_m)t(FAF^{-1}+\del(F)F^{-1})=B.
		\end{equation}		
		Replacing $\l B$ with the $k$-subalgebra of $k'$ generated by $\l B$, the coefficients of the entries of $B$ and $a_1,\ldots,a_m,a_1^{-1},\ldots,a_m^{-1}$, identity (\ref{eq: for regular}) becomes an identity in $\l B[t]_f$. Thus, for any $c \in \Hom_k(\l B, k)$, we can apply the morphism $c\colon \l B[t]_f\to k(t)$ of differential $k$-algebras to (\ref{eq: for regular}) to obtain
		$$
		(t-a_1^c)\ldots(t-a_m^c)t(F^cA^c(F^c)^{-1}+\del(F^c)(F^c)^{-1})=B^c,
		$$
		an identity in $k(t)$. As $B^c\in k[t]^{n\times n}$ and $a_1^c,\ldots,a_m^c$ are non-zero, this shows that $\del(y)=A^cy$ is regular singular at $0$.	
\end{proof}	


%1 - if A already has singularities in k, it's fine (and should be general enough for anything we want : if we want specific singularities in k, we know that we can find a RS equation with the desired singularities over C(z))
%2 - if we want to keep the statement as is, we run into a problem : singularities could merge. Restricting ourselves to a Zariski open should solve the problem and produce a more general (kinda) lemma.

%Let $A\in k(z)^{n\times n}$. Recall that a pole of an entry of $A$ need not be a singularity of $\del(y)=Ay$. Such points are sometimes called \emph{apparent singularities}. Note, however, that they are not singularities of $\del(y)=Ay$ according to our terminology. To have an unambiguous notation we make the following definition. A point $p\in\P^1(k)$ is a \emph{pole of $\del(y)=Ay$} if
%\begin{itemize}
%	\item $p\in k$ and $p$ is a pole of one of the entries of $A$ or
%	\item $p=\infty$ and $z=0$ is a pole of one of the entries of $-\frac{1}{z^2}A(\frac{1}{z})$.
%\end{itemize}	
%Thus the singularities of $\del(y)=Ay$ are contained in the poles of $\del(y)=Ay$.
%
%


The following lemma shows that the conclusion of Theorem \ref{theo: main} is preserved under specialization.


\begin{lemma}\label{lem:regsing}	
	Let $k\subseteq k'$ be an inclusion of algebraically closed fields and let $S=\{p_1,\ldots,p_{d+1}\}$ be a subset of $\P^1(k)\subseteq\P^1(k')$ with $d+1$ elements. Furthermore, let $A \in k'(z)^{n \times n}$ be such that $\del(y)=Ay$ has all poles contained in $S$, is Fuchsian at $p_1,\ldots,p_d$ and regular singular at $p_{d+1}$. Then there exists a finitely generated $k$-subalgebra $\l B \subseteq k'$ and a monic polynomial $f \in \l B[z]$ such that $A \in \l B[z]_f^{n \times n}$ and for every $c \in \Hom_k(\l B, k)$ the differential equation $\del (y) =A^c y$ has all poles contained in $S$, is Fuchsian at $p_1,\ldots,p_d$ and regular singular at $p_{d+1}$.	
\end{lemma}

\begin{proof}
%	There is no difficulty in finding a finitely generated $k$-subalgebra $\l B_0$ of $k'$ and a monic $f_0\in\l B[z]$ such that $A \in \l B_0[z]_f^{n \times n}$ (cf. the proof of Lemma \ref{lem: specialize regsing at 0}).
	We first assume that $\infty\notin S$. By assumption,
	$B=(z-p_1)\ldots(z_1-p_d)(z-p_{d+1})^eA\in k'[z]^{n\times n}$ for some integer $e\geq 1$.	
%	As the poles of $\del(y)=Ay$ are contained in $S$, there exist (not necessarily distinct) $a_1,\ldots,a_m\in S$ such that $B=(z-a_1)\ldots(z_1-a_m)A$ has entries in $k'[z]$.
		 Let $\l B_0$ be the $k$\=/subalgebra of $k'$ generated by the coefficients of all entries of $B$ and set $f=(z-p_1)\ldots(z_1-p_d)(z-p_{d+1})^e\in k[z]\subseteq\l B_0[z]$.
	Then $A\in \l B_0[z]_{f}^{n\times n}$ and for any $c \in \Hom_k(\l B_0, k)$ the poles in $k$ of the differential equation $\del (y) =A^c y$ are contained in $S$.
	
	For a non-zero rational function $h\in \l B_0[z]_{f}\subseteq k'(z)$, when written as $h=\frac{h_1}{f^m}$, with $h_1\in\l B_0[z]$, the rational function $-\frac{1}{t^2}h(\frac{1}{t})\in k'(t)$ has a pole at $t=0$ if and only if $\deg(h_1)-m(d+e)+2>0$. Thus, if $-\frac{1}{t^2}h(\frac{1}{t})$ does not have a pole at $t=0$, for any specialization $c\in \Hom_k(\l B_0, k)$, the rational function $-\frac{1}{t^2}h^c(\frac{1}{t})$ does not have a pole at $t=0$. So, as $p=\infty$ is not a pole of $\del(y)=Ay$, 
	$p=\infty$ is also not a pole of $\del(y)=A^cy$ for any $c\in\Hom_k(\l B_0,k)$. In summary, for any $c\in\Hom_k(\l B_0,k)$, the poles of $\del(y)=A^cy$ are contained in $S$.
	
%	
%	For $f_0=(z-a_1)\ldots(z_1-a_m)$ we then have $A\in \l B_0[x]_{f_0}^{n\times n}$.
%	
%	
%	The singularities of $\del(y)=Ay$ are contained in $\{a_1,\ldots,a_m\}\cup\{\infty\}$. Without loss of generality we may assume that $a_1,\ldots,a_r$ are regular points for $\del(y)=Ay$ and that $a_{r+1},\ldots,a_{m}$ are singular points for $\del(y)=Ay$.
%	
	 The differential equation $\del(y)=A(t+p_{d+1})y$ is regular singular. Thus, by Lemma \ref{lem: specialize regsing at 0}, there exists a finitely generated $k$-subalgebra $\l B$ of $k'$ and a monic polynomial $f_1\in \l B[t]$ such that $A(t+p_{d+1})\in \l B[t]_{f_1}^{n\times n}$ and $\del(y)=A(t+p_{d+1})^cy$ is regular singular for all $c\in\Hom_k(\l B,k)$. Enlarging $\l B$ if necessary, we may assume that $\l B_0\subseteq \l B$. Furthermore, replacing $f_1\in \l B[t]$ by $f_1(t)f(t+p_{d+1})\in \l B[t]$, we may assume that $f(t+p_{d+1})$ divides $f_1$ in $\l B[t]$.
	 We claim that $\l B$ and $f\in\l B[z]$ have the required properties. 
	 
	 Let $c\in\Hom_k(\l B,k)$. Note that $\l B_0[t+p_{d+1}]_{f(t+p_{d+1})}\subseteq \l B[t]_{f_1}$ and $h(t+p_{d+1})^c=h^c(t+p_{d+1})$ for $h\in \l B_0[z]_{f}$. So  $A(t+p_{d+1})^c=A^c(t+p_{d+1})$.
%	 	 $$\l B_0[z]_{f_0}\subseteq \l B[z]_{f_0}=\l B[z-p_{d+1}]_{f_0}\subseteq \l B[z-p_{d+1}]_{f_1(z-p_{d+1})}=\l B[z]_f$$ and $A(t+p_{d+1})^c=A^c(t+p_{d+1})$ for every $c\in\Hom_k(\l B,k)$.
	  As $\del(y)=A(t+p_{d+1})^cy$ is regular singular, this shows that $\del (y)=A^cy$ is a regular singular at $p_{d+1}$. As explained in the second paragraph of this proof, the poles of $\del (y)=A^cy$ are contained in $S$.
	 Finally, from $(z-p_1)\ldots(z_1-p_d)(z-p_{d+1})^eA\in \l B[z]$, we obtain $(z-p_1)\ldots(z_1-p_d)(z-p_{d+1})^eA^c\in k[z]$, showing that $\del (y)=A^cy$ is Fuchsian at $p_1,\ldots,p_{d}$.
	 
	 In case $\infty\in S$, a slight modification of the above proof yields the result. In detail, if $\infty$ is among the Fuchsian points, say $p_1=\infty$, then, with $f=(z-p_2)\ldots(z-p_d)(z-p_{d+1})^e\in k[z]$, we can proceed as above, noting that $-\frac{1}{t}A(\frac{1}{t})\in k'[[t]]^{n\times n}$ implies $-\frac{1}{t}A^c(\frac{1}{t})\in k[[t]]^{n\times n}$ for any $c\in\Hom_k(\l B_0,k)$ (cf. the second paragraph of this proof).
	 
	 If $\infty$ is only regular singular, i.e., $p_{d+1}=\infty$, then, with $f=(z-p_1)\ldots(z-p_d)\in k[z]$ we can proceed similarly to the above, but using that, as $\del(y)=-\frac{1}{t^2}A(\frac{1}{t})y$ is regular singular, there exists, by Lemma~\ref{lem: specialize regsing at 0}, a finitely generated $k$-subalgebra $\l B$ of $k'$ and a monic polynomial $f_1\in\l B[t]$ such that $-\frac{1}{t^2}A(\frac{1}{t})\in \l B[t]_{f_1}^{n\times n}$ and $\del(y)=(-\frac{1}{t^2}A(\frac{1}{t}))^cy$ is regular singular for every $c\in\Hom_k(\l B,k)$. Enlarging $\l B$ if necessary, we can assume that $\l B_0\subseteq \l B$.  Writing $f(\frac{1}{t})=\frac{cf_2}{t^m}$ with $c\in k$ non-zero and $f_2\in k[t]$ monic and replacing $f_1$ by $f_1tf_2$, we can assume that $t$ and $f_2$ divide $f_1$ in $\l B[t]$. As $\l B_0[\frac{1}{t}]\subseteq \l B_0[t]_t\subseteq \l B[t]_{f_1}$ and $\frac{1}{f(\frac{1}{t})}=\frac{c^{-1}t^m}{f_2}\in \l B[t]_{f_1}$, we see that $\l B_0[\frac{1}{t}]_{f(\frac{1}{t})}\subseteq \l B[t]_{f_1}$. Furthermore, for $h\in\l B_0[z]_f$, we have $-\frac{1}{t^2}h(\frac{1}{t})\in\l B[t]_{f_1}$ and $\left(-\frac{1}{t^2}h(\frac{1}{t})\right)^c=-\frac{1}{t^2}h^c(t)$ for any $c\in\Hom_k(\l B,k)$. Thus $\left(-\frac{1}{t^2}A(\frac{1}{t})\right)^c=-\frac{1}{t^2}A(\frac{1}{t})^c$ and  $\del(y)=-\frac{1}{t^2}A(\frac{1}{t})^cy$ is regular singular, i.e., $\del(y)=A^cy$ is regular singular at $\infty$ for every $c\in\Hom_k(\l B,k)$. It follows that $\l B$ and $f$ have the required properties.
%	 
%	 we claim that the lemma holds with $\l B=\l B_1$ and $f=zf_0(z)f_2(z)f_1(z)\in\l B[z]$. Indeed, as 
%	 
%	 
%	  Enlarging $\l B_1$ if necessary, we can assume that $\l B_0\subseteq \l B_1$. Then,
%	 
%	  claim of the lemma then holds with $\l B=\l B_1$ and $f=f_0(z)f_1(z)$.
%	 
%	For the point $p=\infty$ we proceed similarly. % If $p=\infty$ is a regular (or regular singular) point for $\del(y)=Ay$
%	As $\del(y)=-\frac{1}{t^2}A(\frac{1}{t})y$ is regular singular, there exists, by Lemma~\ref{lem: specialize regsing at 0}, a finitely generated $k$-subalgebra $\l B_\infty$ of $k'$ and a monic polynomial $f_\infty\in\l B_\infty[t]$ such that $-\frac{1}{t^2}A(\frac{1}{t})\in \l B_\infty[t]_{f_\infty}^{n\times n}$ and $\del(y)=-\frac{1}{t^2}A(\frac{1}{t})^cy$ is regular singular for every $c\in\Hom_k(\l B_\infty,k)$. We may assume that $\l B_0\subseteq \l B_\infty$. Then $-\frac{1}{t^2}A(\frac{1}{t})^c=-\frac{1}{t^2}A^c(\frac{1}{t})$ and we see that $p=\infty$ is a regular singular point for $\del (y)=A^cy$ for all $c\in\Hom_k(\l B_\infty,k)$.
%	
%	Let $\l B$ be the $k$-subalgebra of $k'$ generated by $\l B_1,\ldots,\l B_m,\l B_\infty$ and set $f=f_0f_1\ldots f_mf_\infty\in\l B[z]$. We claim that $\l B$ and $f$ have the desired properties.
%	
%	
%	
%	Clearly $A\in\l B[z]_f^{n\times n}$. Let $c\in\Hom_k(\l B,k)$. As $\l B_0\subseteq \l B$ (so that $c$ restricts to a morphism $\l B_0\to k$), it follows from the first two paragraphs, that the poles of $\del(y)=Ay$ are contained in $S$. Similarly, as $c$ restricts to morphisms on $\l B_1,\ldots,\l B_m, \l B_\infty$, it follows from the above that $a_1,\ldots,a_m,\infty$ are regular singular points for $\del(y)=A^cy$.
%	As the poles of $\del(y)=A^cy$ are contained in $\{a_1,\ldots,a_m,\infty\}$, and all these points are regular singular, it follows that $\del(y)=A^cy$ is regular singular.	
%%	
%%	We claim that $\del(y)=A^cy$ is regular singular with singularities in $S$.
%	Note that $c\in\Hom_k(\l B,k)$ restricts to $c\in\Hom_k(\l B_i,k)$ for $k\in\{1,\ldots,m,\infty\}$. It follows from the above that $a_1^c,\ldots,a_r^c$ are regular points and $a_{r+1}^c,\ldots,a_m^c$ are regular singular points for $\del(y)=A^cy$. Moreover, $p=\infty$ is a regular (or regular singular) point for $\del(y)=A^cy$ if $p=\infty$ is a regular (or regular singular) point for $\del(y)=Ay$.
%	
%	
%	As $(z-a_1^c)\ldots(z-a_m^c)A^c=B^c\in  k[z]^{n\times n}$, the singularities of $\del(y)=A^cy$ are contained in $\{a_1^c,\ldots,a_m^c,\infty\}$. Because $a_1^c,\ldots,a_r^c$ are regular points for $\del(y)=A^cy$, in fact, the singularities of $\del(y)=A^cy$ are contained in $\{a_{r+1}^c,\ldots,a_m^c,\infty\}$. As all these points are regular singular for $\del(y)=A^c y$, we see that $\del(y)=A^c y$ is regular singular.
%	
%	By assumption, the singularities $a_{r+1},\ldots,a_m$ of $\del(y)=Ay$ are in $k$. So $a_i^c=a_i\in S$ for $i=r+1,\ldots,m$. If $\infty \in S$, this shows that the singularities of $\del(y)=A^cy$ are contained in $S$. If $\infty \notin S$, then $p=\infty$ is a regular point for $\del(y)=Ay$ and therefore $p=\infty$ is a regular point for $\del(y)=A^cy$. Thus, also in this case, the singularities of $\del(y)=A^cy$ are contained in $S$.
%	
%	Note that it is possible that a singularity of $\del(y)=Ay$ is a regular point for $\del(y)=A^cy$. This is no problem as we are only claiming that the singularities are contained in $S$.
%%	
%%	
%	
%	First, notice that the singularities of $A^c$ will be a subset of the singularities of $A$ : since they all lie in $\P^1(k)$, applying $c$ will leave them unchanged. Some singularities might disappear (say, if we have a coefficient of the form $a/z$ and $a$ is mapped to $0$) but that's not a problem because it just means we have less singularities to check.\\
%	We need to prove that $A^c$ is regular singular locally at every point, i.e. that $A^c$ is locally gauge equivalent to a matrix with a pole of order at most $1$ at every point. Since $A$ has only finitely many singularities, we prove that every singularity contributes a finite amount of generators to $\l B$ and a monic factor to $f$. We assume via a coordinate change that the singularity is at $0$.\\
%	\smallskip\noindent
%	Since $A$ is regular singular, Lemma \ref{lem:basis} ensures the existence of a $F \in \GL_{n}\big(\C(z)\big)$ and a matrix with regular coefficients $zB$ such that $F^{-1}(z\frac{d}{dz} - zA) F = z\frac{d}{dz} - zB$.\\
%	Moreover, there exists a monic polynomial $f_1 \in \C[z]$ such that $f_1 A$ and $f_1 F$ are in $\C[z]^{n \times n}$, ie $A$ and $F$ are in $\C[z]_{f_1}^{n \times n}$. We add all the coefficients of $f_1$, as well as those of the polynomial coefficients of $f_1A$ and $f_1F$ to $\l B_1$, so $A,F \in \l B_1[z]_{f_1}$\\
%	Next, we want $F$ to be invertible : we replace $f_1$ with $f_1\det(f_1F) \in \l B_1[z]$, such that :
%	\begin{align*}
%		F^{-1} &= \det(F)^{-1} \text{adj}(F) \\
%		&= f_1 \det(f_1F)^{-1} \text{adj}(f_1F)
%	\end{align*}
%	and by adding the inverse of the leading coefficient of this polynomial to $\l B_1$, we can assume it is monic (and will thus not specialize to the zero polynomial).\\
%	We construct the algebra $\l B$ by iterating this process on other singularities.\\
%	Applying the specialization $c$, we get that :
%	$$F^{c}\lp z\frac{d}{dz}-zA^c\rp \big(F^c\big)^{-1} = z\frac{d}{dz} - zB^c.$$
%	Since $zB^c$ has regular coefficients, the matrix $A^c$ is locally gauge equivalent to matrices with poles of order at most $1$, meaning it is regular singular.
\end{proof}

%\noindent
%We now only need a few lemmas to prove our final result.
%
%


%
%\begin{lemma}\label{lem:regsing}	
%	Let $k\subseteq k'$ be an inclusion of algebraically closed fields and let $S$ be a finite subset of $\P^1(k)\subseteq\P^1(k')$. Furthermore, let $A \in k'(z)^{n \times n}$ be such that $\del(y)=Ay$ is regular singular with poles contained in $S$. Then there exists a finitely generated $k$-subalgebra $\l B \subseteq k'$, and a monic polynomial $f \in \l B[z]$ such that $A \in \l B[z]_f^{n \times n}$ and for every $c \in \Hom_k(\l B, k)$ the differential equation $\del (y) =A^c y$ is regular singular with poles contained in $S$.	
%\end{lemma}
%
%\begin{proof}
%	%	There is no difficulty in finding a finitely generated $k$-subalgebra $\l B_0$ of $k'$ and a monic $f_0\in\l B[z]$ such that $A \in \l B_0[z]_f^{n \times n}$ (cf. the proof of Lemma \ref{lem: specialize regsing at 0}).
%	As the poles of $\del(y)=Ay$ are contained in $S$, there exist (not necessarily distinct) $a_1,\ldots,a_m\in S$ such that $B=(z-a_1)\ldots(z_1-a_m)A$ has entries in $k'[z]$. Let $\l B_0$ be the $k$\=/subalgebra of $k'$ generated by the coefficients of all entries of $B$ and set $f_0=(z-a_1)\ldots(z_1-a_m)$. 
%	Then $A\in \l B_0[z]_{f_0}^{n\times n}$ and for any $c \in \Hom_k(\l B_0, k)$ the poles in $k$ of the differential equation $\del (y) =A^c y$ are contained in $\{a_1,\ldots,a_m\}\subseteq S$.
%	
%	For a non-zero rational function $h\in \l B_0[z]_{f_0}\subseteq k'(z)$, when written as $h=\frac{h_1}{f_0^e}$, with $h_1\in\l B_0[z]$, the rational function $-\frac{1}{z^2}h(\frac{1}{z})\in k'(z)$ has a pole at $z=0$ if and only if $\deg(h_1)-em+2>0$. Thus, if $-\frac{1}{z^2}h(\frac{1}{z})$ does not have a pole at $z=0$, for any specialization $c\in \Hom_k(\l B_0, k)$, the rational function $-\frac{1}{z^2}h^c(\frac{1}{z})$ does not have a pole at $z=0$. So if $p=\infty$ is not a pole of $\del(y)=Ay$, then $p=\infty$ is also not a pole of $\del(y)=A^cy$ for any $c\in\Hom_k(\l B_0,k)$. In summary, for any $c\in\Hom_k(\l B_0,k)$, the poles of $\del(y)=A^cy$ are contained in $S$.
%	
%	%	
%	%	For $f_0=(z-a_1)\ldots(z_1-a_m)$ we then have $A\in \l B_0[x]_{f_0}^{n\times n}$.
%	%	
%	%	
%	%	The singularities of $\del(y)=Ay$ are contained in $\{a_1,\ldots,a_m\}\cup\{\infty\}$. Without loss of generality we may assume that $a_1,\ldots,a_r$ are regular points for $\del(y)=Ay$ and that $a_{r+1},\ldots,a_{m}$ are singular points for $\del(y)=Ay$.
%	%	
%	For $i=1,\ldots,m$, the differential equation $\del(y)=A_iy$, where $A_i=A(t+a_i)\in k'(t)^{n\times n}$ is regular singular. Thus, by Lemma \ref{lem: specialize regsing at 0}, there exists a finitely generated $k$-subalgebra $\l B_i$ of $k'$ and a monic polynomial $f_i\in \l B_i[t]$ such that $A_i\in \l B_i[t]_{f_i}^{n\times n}$ and $\del(y)=A_i^cy$ is regular singular for all $c\in\Hom_k(\l B_i,k)$. Enlarging $\l B_i$ if necessary, we may assume that $\l B_0\subseteq \l B_i$. Then $A_i^c=A(t+a_i)^c=A^c(t+a_i)$ for every $c\in\Hom_k(\l B_i,k)$. As $\del(y)=A_i^cy$ is regular singular, this shows that $p=a_i$ is a regular singular point for $\del (y)=A^cy$.
%	
%	For the point $p=\infty$ we proceed similarly. % If $p=\infty$ is a regular (or regular singular) point for $\del(y)=Ay$
%	As $\del(y)=-\frac{1}{t^2}A(\frac{1}{t})y$ is regular singular, there exists, by Lemma~\ref{lem: specialize regsing at 0}, a finitely generated $k$-subalgebra $\l B_\infty$ of $k'$ and a monic polynomial $f_\infty\in\l B_\infty[t]$ such that $-\frac{1}{t^2}A(\frac{1}{t})\in \l B_\infty[t]_{f_\infty}^{n\times n}$ and $\del(y)=-\frac{1}{t^2}A(\frac{1}{t})^cy$ is regular singular for every $c\in\Hom_k(\l B_\infty,k)$. We may assume that $\l B_0\subseteq \l B_\infty$. Then $-\frac{1}{t^2}A(\frac{1}{t})^c=-\frac{1}{t^2}A^c(\frac{1}{t})$ and we see that $p=\infty$ is a regular singular point for $\del (y)=A^cy$ for all $c\in\Hom_k(\l B_\infty,k)$.
%	
%	Let $\l B$ be the $k$-subalgebra of $k'$ generated by $\l B_1,\ldots,\l B_m,\l B_\infty$ and set $f=f_0f_1\ldots f_mf_\infty\in\l B[z]$. We claim that $\l B$ and $f$ have the desired properties.
%	
%	
%	
%	Clearly $A\in\l B[z]_f^{n\times n}$. Let $c\in\Hom_k(\l B,k)$. As $\l B_0\subseteq \l B$ (so that $c$ restricts to a morphism $\l B_0\to k$), it follows from the first two paragraphs, that the poles of $\del(y)=Ay$ are contained in $S$. Similarly, as $c$ restricts to morphisms on $\l B_1,\ldots,\l B_m, \l B_\infty$, it follows from the above that $a_1,\ldots,a_m,\infty$ are regular singular points for $\del(y)=A^cy$.
%	As the poles of $\del(y)=A^cy$ are contained in $\{a_1,\ldots,a_m,\infty\}$, and all these points are regular singular, it follows that $\del(y)=A^cy$ is regular singular.	
%	%	
%	%	We claim that $\del(y)=A^cy$ is regular singular with singularities in $S$.
%	%	Note that $c\in\Hom_k(\l B,k)$ restricts to $c\in\Hom_k(\l B_i,k)$ for $k\in\{1,\ldots,m,\infty\}$. It follows from the above that $a_1^c,\ldots,a_r^c$ are regular points and $a_{r+1}^c,\ldots,a_m^c$ are regular singular points for $\del(y)=A^cy$. Moreover, $p=\infty$ is a regular (or regular singular) point for $\del(y)=A^cy$ if $p=\infty$ is a regular (or regular singular) point for $\del(y)=Ay$.
%	%	
%	%	
%	%	As $(z-a_1^c)\ldots(z-a_m^c)A^c=B^c\in  k[z]^{n\times n}$, the singularities of $\del(y)=A^cy$ are contained in $\{a_1^c,\ldots,a_m^c,\infty\}$. Because $a_1^c,\ldots,a_r^c$ are regular points for $\del(y)=A^cy$, in fact, the singularities of $\del(y)=A^cy$ are contained in $\{a_{r+1}^c,\ldots,a_m^c,\infty\}$. As all these points are regular singular for $\del(y)=A^c y$, we see that $\del(y)=A^c y$ is regular singular.
%	%	
%	%	By assumption, the singularities $a_{r+1},\ldots,a_m$ of $\del(y)=Ay$ are in $k$. So $a_i^c=a_i\in S$ for $i=r+1,\ldots,m$. If $\infty \in S$, this shows that the singularities of $\del(y)=A^cy$ are contained in $S$. If $\infty \notin S$, then $p=\infty$ is a regular point for $\del(y)=Ay$ and therefore $p=\infty$ is a regular point for $\del(y)=A^cy$. Thus, also in this case, the singularities of $\del(y)=A^cy$ are contained in $S$.
%	%	
%	%	Note that it is possible that a singularity of $\del(y)=Ay$ is a regular point for $\del(y)=A^cy$. This is no problem as we are only claiming that the singularities are contained in $S$.
%	%%	
%	%%	
%	%	
%	%	First, notice that the singularities of $A^c$ will be a subset of the singularities of $A$ : since they all lie in $\P^1(k)$, applying $c$ will leave them unchanged. Some singularities might disappear (say, if we have a coefficient of the form $a/z$ and $a$ is mapped to $0$) but that's not a problem because it just means we have less singularities to check.\\
%	%	We need to prove that $A^c$ is regular singular locally at every point, i.e. that $A^c$ is locally gauge equivalent to a matrix with a pole of order at most $1$ at every point. Since $A$ has only finitely many singularities, we prove that every singularity contributes a finite amount of generators to $\l B$ and a monic factor to $f$. We assume via a coordinate change that the singularity is at $0$.\\
%	%	\smallskip\noindent
%	%	Since $A$ is regular singular, Lemma \ref{lem:basis} ensures the existence of a $F \in \GL_{n}\big(\C(z)\big)$ and a matrix with regular coefficients $zB$ such that $F^{-1}(z\frac{d}{dz} - zA) F = z\frac{d}{dz} - zB$.\\
%	%	Moreover, there exists a monic polynomial $f_1 \in \C[z]$ such that $f_1 A$ and $f_1 F$ are in $\C[z]^{n \times n}$, ie $A$ and $F$ are in $\C[z]_{f_1}^{n \times n}$. We add all the coefficients of $f_1$, as well as those of the polynomial coefficients of $f_1A$ and $f_1F$ to $\l B_1$, so $A,F \in \l B_1[z]_{f_1}$\\
%	%	Next, we want $F$ to be invertible : we replace $f_1$ with $f_1\det(f_1F) \in \l B_1[z]$, such that :
%	%	\begin{align*}
%		%		F^{-1} &= \det(F)^{-1} \text{adj}(F) \\
%		%		&= f_1 \det(f_1F)^{-1} \text{adj}(f_1F)
%		%	\end{align*}
%	%	and by adding the inverse of the leading coefficient of this polynomial to $\l B_1$, we can assume it is monic (and will thus not specialize to the zero polynomial).\\
%	%	We construct the algebra $\l B$ by iterating this process on other singularities.\\
%	%	Applying the specialization $c$, we get that :
%	%	$$F^{c}\lp z\frac{d}{dz}-zA^c\rp \big(F^c\big)^{-1} = z\frac{d}{dz} - zB^c.$$
%	%	Since $zB^c$ has regular coefficients, the matrix $A^c$ is locally gauge equivalent to matrices with poles of order at most $1$, meaning it is regular singular.
%\end{proof}



\subsection{Proof of the main result}

We are now prepared to prove our main theorem.



\begin{theorem}\label{thm:inv}
	Let $G$ be a closed subgroup of $\GL_{n,k}$ that can be generated by $d$ elements and let $S=\{p_1,\ldots,p_{d+1}\}$ be a subset of $\P^1(k)$ with $d+1$ elements. Then there exists a matrix $A\in k(z)^{n\times n}$ such that the differential equation $\del(y)=Ay$ has all poles contained in $S$, is Fuchsian at $p_1,\ldots,p_d$, regular singular at $p_{d+1}$ and has differential Galois group $G$.
\end{theorem}



\begin{proof}
	%Fix an algebraic group $G$ over $k$, and a finite $S \subseteq \P^1(k)$.
	 We roughly follow the short solution to the inverse problem from \cite[Section 5.2]{fengwib}. The algebraic group $G$ (together with the closed embedding into $\GL_{n,k}$) descends to a finitely generated field $k_0$, i.e., there exist a subfield $k_0$ of $k$, finitely generated over $\mathbb{Q}$, and a closed subgroup $G_0$ of $\GL_{n,k_0}$ such that $(G_0)_k=G$ (as closed subgroups of $\GL_{n,k}$).
	 
	 Let $g_1,\ldots,g_d\in G(k)$ generate $G$. Enlarging $k_0$ if necessary, we can assume that $g_1,\ldots,g_d\in G_0(k_0)\subseteq G_0(k)=G(k)$ and $S\subseteq \P^1(k_0)$.
	 
	 
	 %It follows from Lemma ?? that $g_1,\ldots,g_d\in G_0(k_0)$ generate $G_0$. 
	 
	 Let $k_1\subseteq k$ be the algebraic closure of $k_0$ and set $G_1=(G_0)_{k_1}$. Then $G_1$ is a closed subgroup of $\GL_{n,k_1}$. %As  that can be generated by $d$ elements (??). 
	 As $k_1$ is countable, there exists an embedding $k_1\hookrightarrow \C$ (that we fix).
	 
	 Because $g_1,\ldots,g_d\in G_0(k_0)\subseteq G_0(k_1)=G_1(k_1)$ generate $G=(G_1)_k$, it follows from Lemma~\ref{lem:groupgen} that $g_1,\ldots,g_d$ generate $G_1$. The same lemma, but applied in the reverse direction, shows that $(G_1)_{\C}$ is generated by $g_1,\ldots,g_d$. In particular, $(G_1)_{\C}$ is a closed subgroup of $\GL_{n,\C}$ that can be generated by $d$ elements.
	 
	 As noted in the introduction, Theorem \ref{thm:inv} holds for $k=\C$ as a consequence of the (weak) solution of the Riemann-Hilbert problem (see \cite[Theorem 3.2.1]{AnasovBolibruch:TheRiemannHilbertProblem}, \cite[Theorem 12.4]{Saul} or \cite[Part I, Theorem 4.49]{MitschiSauzin}) and Schlesinger's density theorem (see e.g., \cite[Part I, Theorem 2.28]{MitschiSauzin}).
	 	 %\cite[Theorem 5.8]{SVDP}). 
%	 In fact, the (weak) solution of the Riemann-Hilbert problem yields a slightly stronger statement than Theorem \ref{thm:inv} in the case $k=\mathbb{C}$: The linear differential equation can be chosen such that its poles are contained in $S$. (See e.g., \cite[Part I, Theorem 4.49]{MitschiSauzin} or \cite[Theorem 12.4]{Saul}.)	 
	 Thus, there exists a matrix $A\in \mathbb{C}(z)^{n\times n}$ such that $\del(y)=Ay$ has poles contained in $S$, is Fuchsian at $p_1,\ldots,p_d$, regular singular at $p_{d+1}$ and has differential Galois group $(G_1)_{\mathbb{C}}$. 
	 
%	 So applying the theorem to $(G_1)_{\C}\leq\GL_{n,\C}$ and $S\subseteq\P^1(k_0)\subseteq\P^1(\C)$ yields a matrix $\mathcal{A}\in\C(z)^{n\times n}$ such that $\del(y)=\mathcal{A}y$ is regular singular with singularities in $S$ and has differential Galois group $(G_1)_{\C}$. 
%	 
%	 algebraic group $G'$ over $k_1$ such that $G'_k = G$. By extending $k_0$, we can ensure the generators of $G(k)$ are already in $G'(k)$, and thus $G'$ is generated by $d$ elements. We may also ensure that $S \subseteq \P^1(k_0)$. Since $k_0$ is finitely generated, both it and $k_1$ can be embedded into $\C$. We fix such an embedding. Notice that because of Lemma \ref{lem:groupgen}, $G'_\C$ is also generated by $d$ elements.
%	
%	\newpar
%	There exists a regular singular matrix $A$ with singularities in $S$ (seen as a subset of $\P^1(k_1) \subseteq \P^1(\C)$) such that $\partial y=By$ has Galois group $G'_\C$, because every algebraic group over $\C$ generated by $d$ elements is the Galois group of some regular singular differential equation over $\C(z)$ with singularities in $S$ (\cite{SVDP}, Theorem 5.12).
%	

	Applying Lemma \ref{lem:regsing} to the inclusion $k_1\subseteq \C$ and the differential equation $\del(y)=Ay$ yields a finitely generated $k_1$-subalgebra $\l B_0$ of $\C$ and a monic $f_0\in \l B_0[z]$ with $A\in \l B_0[z]_{f_0}^{n\times n}$ such that $\del(y)=A^cy$ has all poles contained in $S$, is Fuchsian at $p_1,\ldots,p_d$ and regular singular at $p_{d+1}$ for all $c\in\Hom_{k_1}(\l B_0,k_1)$. By Lemma~\ref{lem:spread}, there exists a finitely generated $k_1$-subalgebra $\l B$ of $\mathbb{C}$ containing $\l B_0$, a monic polynomial $f \in \l B[z]$ such that $f_0$ divides $f$, and a differential $(G_1)_{\l B}$\=/torsor $\l R/\l B[z]_f$ for $\del(y)=Ay$ such that $\l R$ is flat over $\l B[z]_f$ and $\l R\otimes_{\l B[z]_f}K(z)$ is $\del$-simple, where $K$ is the algebraic closure of the field of fractions of $\l B$.
	By Theorem \ref{thm:spe}, there exists a $c \in \Hom_{k_1}(\l B, k_1)$ such that $R^c = \l R \otimes_{\l B[z]_f} k_1(z)$ is a Picard-Vessiot ring for $\del(y)=A^cy$ (over $k_1(z)$) with differential Galois group $((G_1)_{\l B})_{k_1}=G_1$.
	By construction of $\l B_0$ and $f_0$, and using that $A\in \l B_0[z]_{f_0}^{n\times n}\subseteq \l B[z]_f^{n\times n}$, we see that $\del(y)=A^cy$ has all poles contained in $S$, is Fuchsian at $p_1,\ldots,p_d$ and regular singular at $p_{d+1}$.
	
	Finally, by Lemma \ref{lem:basechange}, the equation $\del(y)=A^cy$, considered as a differential equation over $k(z)$, has the required properties, i.e., $\del(y)=A^cy$ (over $k(z)$) has all poles contained in $S$, is Fuchsian at $p_1,\ldots,p_d$, regular singular at $p_{d+1}$ and has differential Galois group $(G_1)_k=G$.	
	% Let $\l B_0$ and $f_0$ be the $k_1$-algebra and monic polynomial given by Lemma \ref{lem:regsing} for $A$.	
	%	\vspace{40mm}
% The Picard-Vessiot ring of the equation can then be described using Lemma \ref{lem:spread} : there exists
%	
%	\begin{itemize}
%		\ibul a finitely generated $k_1$-algebra $\l B$ containing $\l B_0$,
%		\ibul a monic polynomial $f \in \l B[z]$ such that $f_0 | f$ and $A \in \l B[z]_f$,
%		\ibul a $G'_{\l B}$-torsor $\l R/\l B[z]_f$ such that $R^{\text{gen}}$ is differentially simple,
%		\ibul and a fundamental matrix $\l Y \in \GL_n(\l R)$ with $\frac{d}{dz}\l Y = A \l Y$.
%	\end{itemize}
%	By Theorem \ref{thm:spe} there exists a $c \in \Hom_{k_1}(\l B, k_1)$ such that $R^c = \l R \otimes_{\l B[z]_f} k(z)$ is a Picard-Vessiot ring for $\del(y)=A^cy$ (over $k_1(z)$). Since $\l R$ is generated by the entries of $\l Y$, $R^c$ is generated by the entries of $Y^c := \l Y \otimes 1$ (which are exactly the entries of $\l Y$ after applying $c$), which satisfies the equation $\frac{d}{dz} Y^c = B^c Y^c$, so $R^c$ is a Picard-Vessiot ring for this equation. Set $A := B^c$.\\
%	Its Galois group will be $\big(G'_{\l B}\big)_{k_1} = G'$, showing that $G'$ will be the Galois group over $k_1(z)$ of a regular singular matrix equation.
\end{proof}

%The above proof shows that the matrix $A$ in Theorem \ref{thm:inv} can be chosen such that the poles of $\del(y)=Ay$ are contained in $S$.  ??

%\noindent
%Since every algebraic group is finitely generated (\cite{SVDP}, Lemma 5.13), we have the following corollary :

	\begin{corollary}\label{cor:inv}
		Any algebraic group over $k$ is the differential Galois group of a regular singular differential equation $\del (y) = A y$ over $k(z)$.
	\end{corollary}
	\begin{proof}
	As any (affine) algebraic group is isomorphic to a closed subgroup of some general linear group (\cite[Cor. 4.10]{Milne:AlgebraicGroups}) and every algebraic group (in characteristic zero) is finitely generated (\cite{SVDP}, Lemma 5.13), this follows from Theorem \ref{thm:inv}.
	\end{proof}	




\subsection{An interpretation in terms of proalgebraic groups}


In this final section we offer an interpretation of our main result (Theorem \ref{thm:inv}) in terms of proalgebraic groups.


Recall that the \emph{free proalgebraic group} $\Gamma_d$ on $d$ generators (over $k$) is characterized by the following universal property of the map $\iota\colon X\to \Gamma_d(k)$ from a $d$-element set $X$ into the $k$-points of $\Gamma_d$: For every proalgebraic group $G$ and every map $\varphi\colon X\to G(k)$ there exists a unique morphism $\phi\colon \Gamma_d\to G$ of proalgebraic groups such that
$$
\begin{tikzcd}  X \arrow[d, "\varphi"'] \arrow{r}{\iota} & \Gamma_d(k) \arrow[ld, dashrightarrow, "\phi"]  \\  G(k)  \end{tikzcd}
$$
commutes. The proalgebraic group $\Gamma_d$ can be constructed as the fundamental group of the neutral tannakian category of all (finite dimensional, $k$-linear) representations of the (abstract) free group $F_d$ on $d$ generators. In other words, $\Gamma_d$ is the proalgebraic completion (or proalgebraic hull) of $F_d$. (See \cite{proalg} for more background on free proalgebraic groups.) 


Let $S$ be a subset of $\P^1(k)$ with $d+1$ elements and let $\Gamma_S$ be the differential Galois group of the family of all regular singular differential equations over $k(z)$ with singularities contained in $S$.
For $k=\C$, it is an immediate corollary of the Riemann-Hilbert correspondence (\cite[Theorem~6.15]{SVDP}) that $\Gamma_S$ is isomorphic to $\Gamma_d$. It therefore seems natural to expect that $\Gamma_S$ is isomorphic to $\Gamma_d$ in general. This question was already raised in \cite[Section~4.1]{Wibmer:RegSing}. While we are far from being able to show that indeed $\Gamma_S$ is isomorphic to $\Gamma_d$, Theorem \ref{thm:inv} can be seen as a small step into the right direction. It implies that every algebraic quotient of $\Gamma_d$ is also a quotient of $\Gamma_S$.



%\noindent
%We can formulate this result in terms of free proalgebraic groups (see \cite{proalg} for a reference on free proalgebraic groups in general) :


\begin{corollary}\label{cor:proalg}
	Let $G$ be an algebraic group over $k$. If $G$ is a quotient of $\Gamma_d$, then $G$ is also a quotient of $\Gamma_S$.
\end{corollary}
\begin{proof}
If $G$ is a quotient of $\Gamma_d$, then $G$ can be generated by $d$ elements (\cite[Lemma 2.16]{proalg}). After embedding $G$ into some $\GL_{n,k}$, we can apply Theorem \ref{thm:inv} to find a regular singular differential equation $\del(y)=Ay$ with differential Galois group $G$. The Picard-Vessiot ring $R$ of $\del(y)=Ay$ is canonically embedded in the Picard-Vessiot ring $R_S$ of the family of all regular singular differential equations with singularities in $S$. By the second fundamental theorem of differential Galois theory (see e.g., \cite[Theroem 2.11]{AmanoMasuokaTakeuchi}), the differential Galois group of $R$ is a quotient of the differential Galois group of $R_S$, i.e., $G$ is a quotient of $\Gamma_S$ as desired.
\end{proof}	


%\begin{corollary}[Galois group of all regular singular equations]\label{cor:proalg}
%		Let $k$ be an algebraically closed field of characteristic zero, and $S \subset \P^1(k)$ a subset of cardinality $d+1$. All quotients of the free proalgebraic group on $d$ elements, i.e. all algebraic groups generated by $\le d$ elements occur as algebraic quotients of the Galois group of all regular singular equations with singularities in $S$.
%\end{corollary}
%
%\begin{proof}
%Call $\hat{R}_S$ the Picard-Vessiot ring of all regular singular equations with singularities in $S$. The algebraic quotients of $\Gal(\hat{R}_S/k(z))$ are the Galois groups of Picard-Vessiot rings $k(z) \subseteq R \subseteq \hat{R}_S$ such that $(R \otimes_{k(z)}R)^\del$ is finitely generated as a $k$-algebra, or, equivalently, that $R$ is finitely generated as a $k(z)$-algebra, meaning $R$ is a Picard-Vessiot ring for a matrix differential equation.
%\end{proof}
%
%\noindent
%Theorem \ref{cor:proalg} is a nice step in the direction of a description of the Galois group of all regular singular differential equations with singularities in (finite) $S \subseteq \P^1(k)$ over $k(z)$ with algebraically closed $k$ of characteristic zero. The next natural question to ask is whether this group is actually isomorphic to the free proalgebraic group on $d = |S|-1$ elements : this is known (even for infinite $S$) when $k=\C$, but not for any other algebraically closed field.


\printbibliography



Thomas Serafini, Sorbonne Université, IMJ-PRG, 4 place Jussieu, 
75005 Paris, France, \texttt{tserafini@dma.ens.fr} 

\medskip

Michael Wibmer, School of Mathematics, University of Leeds, LS2 9JT, Leeds, United Kingdom, \texttt{m.wibmer@leeds.ac.uk}


\end{document}

